% Options for packages loaded elsewhere
\PassOptionsToPackage{unicode}{hyperref}
\PassOptionsToPackage{hyphens}{url}
\PassOptionsToPackage{dvipsnames,svgnames,x11names}{xcolor}
%
\documentclass[
  letterpaper,
  DIV=11,
  numbers=noendperiod]{scrartcl}

\usepackage{amsmath,amssymb}
\usepackage{iftex}
\ifPDFTeX
  \usepackage[T1]{fontenc}
  \usepackage[utf8]{inputenc}
  \usepackage{textcomp} % provide euro and other symbols
\else % if luatex or xetex
  \usepackage{unicode-math}
  \defaultfontfeatures{Scale=MatchLowercase}
  \defaultfontfeatures[\rmfamily]{Ligatures=TeX,Scale=1}
\fi
\usepackage{lmodern}
\ifPDFTeX\else  
    % xetex/luatex font selection
\fi
% Use upquote if available, for straight quotes in verbatim environments
\IfFileExists{upquote.sty}{\usepackage{upquote}}{}
\IfFileExists{microtype.sty}{% use microtype if available
  \usepackage[]{microtype}
  \UseMicrotypeSet[protrusion]{basicmath} % disable protrusion for tt fonts
}{}
\makeatletter
\@ifundefined{KOMAClassName}{% if non-KOMA class
  \IfFileExists{parskip.sty}{%
    \usepackage{parskip}
  }{% else
    \setlength{\parindent}{0pt}
    \setlength{\parskip}{6pt plus 2pt minus 1pt}}
}{% if KOMA class
  \KOMAoptions{parskip=half}}
\makeatother
\usepackage{xcolor}
\setlength{\emergencystretch}{3em} % prevent overfull lines
\setcounter{secnumdepth}{-\maxdimen} % remove section numbering
% Make \paragraph and \subparagraph free-standing
\makeatletter
\ifx\paragraph\undefined\else
  \let\oldparagraph\paragraph
  \renewcommand{\paragraph}{
    \@ifstar
      \xxxParagraphStar
      \xxxParagraphNoStar
  }
  \newcommand{\xxxParagraphStar}[1]{\oldparagraph*{#1}\mbox{}}
  \newcommand{\xxxParagraphNoStar}[1]{\oldparagraph{#1}\mbox{}}
\fi
\ifx\subparagraph\undefined\else
  \let\oldsubparagraph\subparagraph
  \renewcommand{\subparagraph}{
    \@ifstar
      \xxxSubParagraphStar
      \xxxSubParagraphNoStar
  }
  \newcommand{\xxxSubParagraphStar}[1]{\oldsubparagraph*{#1}\mbox{}}
  \newcommand{\xxxSubParagraphNoStar}[1]{\oldsubparagraph{#1}\mbox{}}
\fi
\makeatother


\providecommand{\tightlist}{%
  \setlength{\itemsep}{0pt}\setlength{\parskip}{0pt}}\usepackage{longtable,booktabs,array}
\usepackage{calc} % for calculating minipage widths
% Correct order of tables after \paragraph or \subparagraph
\usepackage{etoolbox}
\makeatletter
\patchcmd\longtable{\par}{\if@noskipsec\mbox{}\fi\par}{}{}
\makeatother
% Allow footnotes in longtable head/foot
\IfFileExists{footnotehyper.sty}{\usepackage{footnotehyper}}{\usepackage{footnote}}
\makesavenoteenv{longtable}
\usepackage{graphicx}
\makeatletter
\newsavebox\pandoc@box
\newcommand*\pandocbounded[1]{% scales image to fit in text height/width
  \sbox\pandoc@box{#1}%
  \Gscale@div\@tempa{\textheight}{\dimexpr\ht\pandoc@box+\dp\pandoc@box\relax}%
  \Gscale@div\@tempb{\linewidth}{\wd\pandoc@box}%
  \ifdim\@tempb\p@<\@tempa\p@\let\@tempa\@tempb\fi% select the smaller of both
  \ifdim\@tempa\p@<\p@\scalebox{\@tempa}{\usebox\pandoc@box}%
  \else\usebox{\pandoc@box}%
  \fi%
}
% Set default figure placement to htbp
\def\fps@figure{htbp}
\makeatother

\KOMAoption{captions}{tableheading}
\makeatletter
\@ifpackageloaded{tcolorbox}{}{\usepackage[skins,breakable]{tcolorbox}}
\@ifpackageloaded{fontawesome5}{}{\usepackage{fontawesome5}}
\definecolor{quarto-callout-color}{HTML}{909090}
\definecolor{quarto-callout-note-color}{HTML}{0758E5}
\definecolor{quarto-callout-important-color}{HTML}{CC1914}
\definecolor{quarto-callout-warning-color}{HTML}{EB9113}
\definecolor{quarto-callout-tip-color}{HTML}{00A047}
\definecolor{quarto-callout-caution-color}{HTML}{FC5300}
\definecolor{quarto-callout-color-frame}{HTML}{acacac}
\definecolor{quarto-callout-note-color-frame}{HTML}{4582ec}
\definecolor{quarto-callout-important-color-frame}{HTML}{d9534f}
\definecolor{quarto-callout-warning-color-frame}{HTML}{f0ad4e}
\definecolor{quarto-callout-tip-color-frame}{HTML}{02b875}
\definecolor{quarto-callout-caution-color-frame}{HTML}{fd7e14}
\makeatother
\makeatletter
\@ifpackageloaded{caption}{}{\usepackage{caption}}
\AtBeginDocument{%
\ifdefined\contentsname
  \renewcommand*\contentsname{Table of contents}
\else
  \newcommand\contentsname{Table of contents}
\fi
\ifdefined\listfigurename
  \renewcommand*\listfigurename{List of Figures}
\else
  \newcommand\listfigurename{List of Figures}
\fi
\ifdefined\listtablename
  \renewcommand*\listtablename{List of Tables}
\else
  \newcommand\listtablename{List of Tables}
\fi
\ifdefined\figurename
  \renewcommand*\figurename{Figure}
\else
  \newcommand\figurename{Figure}
\fi
\ifdefined\tablename
  \renewcommand*\tablename{Table}
\else
  \newcommand\tablename{Table}
\fi
}
\@ifpackageloaded{float}{}{\usepackage{float}}
\floatstyle{ruled}
\@ifundefined{c@chapter}{\newfloat{codelisting}{h}{lop}}{\newfloat{codelisting}{h}{lop}[chapter]}
\floatname{codelisting}{Listing}
\newcommand*\listoflistings{\listof{codelisting}{List of Listings}}
\makeatother
\makeatletter
\makeatother
\makeatletter
\@ifpackageloaded{caption}{}{\usepackage{caption}}
\@ifpackageloaded{subcaption}{}{\usepackage{subcaption}}
\makeatother

\usepackage{bookmark}

\IfFileExists{xurl.sty}{\usepackage{xurl}}{} % add URL line breaks if available
\urlstyle{same} % disable monospaced font for URLs
\hypersetup{
  colorlinks=true,
  linkcolor={blue},
  filecolor={Maroon},
  citecolor={Blue},
  urlcolor={Blue},
  pdfcreator={LaTeX via pandoc}}


\author{}
\date{}

\begin{document}


\subsection{Nachhaltigkeit?}\label{nachhaltigkeit}

Woher kommen das Konzept, das Leitbild und die Handlungsprinzipien
nachhaltiger Entwicklung? Wie ist das Konzept entstanden, wie hat es
sich entwickelt -- und warum? Die folgenden Ausführungen geben einen
Überblick über die Geschichte der Nachhaltigkeit, ihre politischen
Hintergründe sowie über wichtige Konferenzen und Dokumente. Kurz gesagt:
Es geht um das, was Nachhaltigkeit zu dem gemacht hat, was sie heute
ist. Es geht darum, den Wald hinter den Bäumen zu sehen. Denn nur wer
die Vergangenheit kennt, kann die Zukunft einschätzen.

Nachhaltigkeit ist in den letzten Jahren zu einem Schlagwort geworden --
in der Wissenschaft, der Politik, der Wirtschaft und den Medien. Der
Begriff ist zunächst positiv besetzt, da er mit Langfristigkeit und
Beständigkeit assoziiert wird. Heute ist vieles nachhaltig: der Kaffee,
die Firmenphilosophie, selbst die Thunfischpizza. Dies sind die
Ansprüche einer Konsumgesellschaft mit einem rücksichtslosen Appetit auf
Überfluss. Gleich am Morgen entfernt uns das Shampoo «nachhaltig» die
Schuppen. Dann fahren wir zur Arbeit in eine Firma, die sich mit einer
«nachhaltigen Unternehmensphilosophie» schmückt, und beim Mittagessen
diskutieren wir über nachhaltige Geldanlagen. Wieder zu Hause schieben
wir eine Thunfischpizza in den Ofen -- selbstverständlich aus
«nachhaltiger Erzeugung», wie es auf der Verpackung steht.

Heute begegnet uns der Begriff «Nachhaltigkeit» überall: in Zusammenhang
mit Energie, Mobilität, Gebäudesanierung, Ernährung,
Bevölkerungsentwicklung, betrieblichem Umweltmanagement und Klimaschutz.
Aber auch in Kunst, Kultur, Design und Werbung.

\subsection{Ein normatives Konzept}\label{ein-normatives-konzept}

Um Nachhaltigkeit zu erreichen, benötigen wir nachhaltige Entwicklung:
einen strategischen Prozess, der Veränderungen in unseren
sozial‑ökologischen Systemen und Institutionen erfordert. „Entwicklung``
impliziert eine gesteuerte Verbesserung, und die ethische Frage, was
genau „besser`` bedeutet, ist dabei zentral. Die Verwendung des Begriffs
„Entwicklung`` wird von einigen Autor:innen (z.B. Lang et al.~2014;
Kothari et al.~2019) kritisiert, da er häufig mit ungezügeltem Wachstum
gleichgesetzt wird. Ein ungebremstes Wachstum des „ökologischen
Fussabdrucks`` -- also des Anteils der Biosphäre, der für menschliche
Produktion, Konsum und Abfall genutzt wird -- ist langfristig nicht
nachhaltig. Dennoch gibt es unterschiedliche Interpretationen von
„Entwicklung``, die von wirtschaftlichem Wachstum bis hin zur
Verbesserung der Lebensqualität reichen. Nachhaltige Entwicklung bleibt
daher ein anregendes und umstrittenes Konzept. Nachhaltigkeit und
nachhaltige Entwicklung spielen heute eine Schlüsselrolle in politischen
Diskursen. Der Begriff „Nicht‑Nachhaltigkeit`` verweist auf Zustände
oder Entwicklungen, die als negativ bewertet werden, während
„Nachhaltigkeit`` einen positiven Zustand beschreibt.

Das Konzept nachhaltiger Entwicklung verpflichtet uns zu bestimmten
Werten und Normen, die ethische Ziele und Handlungsregeln festlegen.
Diese Werte und Normen prägen unsere Vorstellungen davon, was wir als
wünschenswerten oder „positiven`` Wandel ansehen. Ein neutraler
Standpunkt ist dabei nicht möglich, da unsere Wahrnehmung durch eine
Vielzahl von Einflüssen geformt wird. Unsere persönlichen Erfahrungen,
unsere Erziehung, unser soziales Umfeld und unsere kulturellen
Hintergründe beeinflussen, wie wir die Welt sehen und interpretieren --
und damit auch, was wir als „sein sollend`` wahrnehmen und welche
Handlungen wir als notwendig ansehen. Ethik spielt eine Schlüsselrolle
dabei, wie die aktuelle Situation verbessert werden kann, indem sie
normative Ziele und Grenzen definiert. Nachhaltige Entwicklung ist somit
ein konzeptioneller Rahmen, der stark von ethischen Überlegungen
getragen wird.

Nachhaltigkeitsherausforderungen wie Armutsbekämpfung und Klimawandel
sind daher ethische Herausforderungen. Die Identifikation von
Situationen als Nachhaltigkeitsprobleme (und damit als „negativ``) sowie
die Auswahl von Lösungen beruhen auf ethischen Wertvorstellungen.
Nachhaltigkeitsziele basieren ebenfalls auf solchen Werten, ebenso wie
auf Wissen über Ursache‑Wirkungs‑Zusammenhänge. Fragen der
Nachhaltigkeit werden häufig als „vertrackte Probleme`` (wicked
problems) bezeichnet, da sie komplex und pluralistisch sind. Die globale
Erwärmung wird sogar als „super wicked problem`` beschrieben, da
Lösungen dringend sind, politische Institutionen häufig unzureichend
agieren und Entscheidungsprozesse unter Kurzsichtigkeit leiden.

\subsection{Globale Herausforderungen als vertrackte
Probleme}\label{globale-herausforderungen-als-vertrackte-probleme}

Mike Hulme, Autor von Why We Disagree About Climate Change (2009),
schlägt vor, den Klimawandel als ein „vertracktes Problem`` (wicked
problem) von enormem Ausmaß und wahrscheinlicher Langlebigkeit zu
betrachten. Dieser Ansatz hilft uns, den Klimawandel nicht als ein
Problem zu sehen, das gelöst werden muss, sondern als eine Bedingung, in
der wir unmittelbar eingebunden sind. Die Kategorisierung von Problemen
als „vertrackt``, also solcher, die sich keiner eindeutigen Lösung
fügen, geht auf die Planungstheoretiker Horst Rittel und Melvin Webber
(1973) zurück. Sie argumentierten, dass Planende mit unvorhersehbarem
menschlichem Verhalten umgehen müssen und dass einige Probleme zu
komplex sind, um vollständig gelöst zu werden. Manche „vertrackten
Probleme`` lassen sich möglicherweise nie endgültig lösen, doch wir
können lernen, besser mit ihnen umzugehen und uns nicht von ihnen
beherrschen zu lassen.

\begin{tcolorbox}[enhanced jigsaw, breakable, left=2mm, colback=white, leftrule=.75mm, opacityback=0, bottomrule=.15mm, toprule=.15mm, arc=.35mm, rightrule=.15mm, colframe=quarto-callout-note-color-frame]

\vspace{-3mm}\textbf{Vertrackte Probleme}\vspace{3mm}

Laut Bannink und Trommel (2019) entstehen vertrackte Probleme an der
Schnittstelle zwischen faktischer Ungewissheit und einer Heterogenität
von Präferenzen und Interessen. Vertrackte Probleme sind gekennzeichnet
durch (Rittel \& Weber, 1973; Alford, 2017; Sediri, 2020):

\begin{itemize}
\item
  Komplexität: Vertrackte Probleme sind durch zahlreiche miteinander
  verknüpfte Faktoren und Wechselwirkungen gekennzeichnet. Es gibt keine
  klaren Ursache‑Wirkungs‑Beziehungen, und Veränderungen in einem
  Bereich können unvorhergesehene Auswirkungen in anderen Bereichen
  haben.
\item
  Normative Konflikte: Vertrackte Probleme werden von unterschiedlichen
  Akteur:innen und Interessengruppen unterschiedlich wahrgenommen und
  interpretiert. Es gibt keine eindeutige Definition oder einen Konsens
  darüber, was genau das Problem ist oder wie es gelöst werden sollte.
\item
  Interdisziplinarität: Vertrackte Probleme erfordern einen
  interdisziplinären Ansatz, da sie verschiedene Themenfelder,
  Perspektiven und Akteur:innen betreffen. Die Lösungsfindung erfordert
  oft die Zusammenarbeit und Koordination verschiedener Disziplinen und
  Fachleute.
\item
  Ungewissheit: Fehlerhafte, unvollständige oder nicht zugängliche
  Informationen über die Problemsituation und über die Stabilität der
  zugrunde liegenden Werte der beteiligten Variablen erschweren das
  Verständnis.
\item
  Keine endgültige Lösung: Vertrackte Probleme entziehen sich
  eindeutigen Lösungen. Sie sind dynamisch, verändern sich im Laufe der
  Zeit und erfordern kontinuierliche Anpassung und wiederholte
  Bearbeitung.
\end{itemize}

Beispiele für vertrackte Probleme sind u.a. der Klimawandel, die
Armutsbekämpfung, die globale Gesundheitsversorgung und die
Nachhaltigkeit. Die Komplexität und das Zusammenspiel verschiedener
Faktoren in diesen Bereichen machen es schwierig, einfache und
eindeutige Lösungen zu finden. Der Umgang mit vertrackten Problemen
erfordert ein hohes Maß an Reflexion, Zusammenarbeit und den Einsatz von
Methoden des systemischen Denkens.

\end{tcolorbox}

Wie haben wir uns in die heutige Lage manövriert, in der vertrackte
Probleme eine so grosse Herausforderung für eine nachhaltige Entwicklung
darstellen? In diesem Lehrbuch werden wir drei zentrale globale
Herausforderungen analysieren und verstehen lernen:

\begin{itemize}
\item
  die Entstehung der vom Menschen verursachten globalen Erwärmung;
\item
  die Persistenz globaler Armut und steigender Ungleichheiten;
\item
  sowie die Bedrohung durch Artensterben und die Übernutzung natürlicher
  Ressourcen.
\end{itemize}

Die Debatte über den menschlichen Einfluss auf die globale Erwärmung und
den Klimawandel hat sich seit den Veröffentlichungen des
Intergovernmental Panel on Climate Change (IPCC) intensiviert. Die
treibende Kraft hinter der globalen Erwärmung ist unsere Abhängigkeit
von Erdöl und anderen fossilen Brennstoffen für Transport,
(Strom-)Erzeugung, Landwirtschaft und viele Alltagsprodukte. Ein
weiteres vertracktes Problem ist die globale Armut. Obwohl es seit
Jahrzehnten weltweite Bemühungen gibt, die Armut zu verringern, führt
die globalisierte Marktwirtschaft in bestimmten Regionen häufig zu
zunehmender Armut, und die Kluft zwischen Arm und Reich scheint
insgesamt größer zu werden (Chancel et al., 2021). Hocheinkommensländer
haben ihren Wohlstand und Lebensstandard auf Kosten anderer
aufrechterhalten. Diese „externalisierten Kosten`` umfassen auch
Umweltzerstörung. So sind der Wohlstand und die Lebensstandards von
Hocheinkommensländern maßgeblich für die Bedrohung durch Artensterben
und die Übernutzung natürlicher Ressourcen verantwortlich (Lessenich,
2018).

Zusätzlich zu den oben genannten „grossen Drei`` -- Klimawandel, Armut
und Biodiversitätsverlust -- gibt es viele weitere globale
Nachhaltigkeitsherausforderungen, wie Entwaldung, Wüstenbildung,
abnehmende Bodenfruchtbarkeit, Rückgang der Fischbestände,
Umweltverschmutzung, Kriege, Konflikte und zunehmende Migration. Für die
Analyse dieser umfassenden Herausforderungen müssen wir auch die
zeitliche Dimension berücksichtigen. Im Jahr 2004 wurden erstmals
Grafiken veröffentlicht, die sozioökonomische und Erdsystemtrends von
1750 bis 2000 darstellen. Sie zeigen einen dramatischen Anstieg und
wurden treffend als „Die grosse Beschleunigung`` (The Great
Acceleration) bezeichnet (Steffen et al., 2004).

Die grosse Beschleunigung beschreibt die exponentielle Wachstumsdynamik
nach der Industriellen Revolution, die zu einem sprunghaften Anstieg der
Produktivität in der Produktionsweise und zu einem erheblichen Zuwachs
an materiellem Wohlstand führte. Gleichzeitig, aber deutlich weniger
beachtet, kam es zu einem massiven Anstieg des Verbrauchs natürlicher
Ressourcen und der Emissionen. Sozioökonomisches Wachstum ging Hand in
Hand mit der Beschleunigung biophysischer Trends. Abbildungen 6 und 7
zeigen Beispiele wichtiger biophysischer und sozioökonomischer
Indikatoren, die alle mit der Industriellen Revolution zu steigen
beginnen. Ab der Mitte des 20. Jahrhunderts wird die Tendenz zu
exponentiellem Wachstum offensichtlich.

\begin{tcolorbox}[enhanced jigsaw, breakable, left=2mm, colback=white, leftrule=.75mm, opacityback=0, bottomrule=.15mm, toprule=.15mm, arc=.35mm, rightrule=.15mm, colframe=quarto-callout-tip-color-frame]

\vspace{-3mm}\textbf{Eine Veranschaulichung exponentiellen Wachstums}\vspace{3mm}

Das Weizenkorn‑ und Schachbrett‑Problem ist ein mathematisches
Gedankenexperiment, das häufig verwendet wird, um das Konzept
exponentiellen Wachstums zu verdeutlichen. In der Geschichte bittet ein
Diener den König, jedes Feld eines Schachbretts mit Weizenkörnern zu
belegen und die Anzahl auf jedem Feld jeweils zu verdoppeln. Das
Wachstum beginnt langsam, doch mit jeder Verdopplung steigt die Anzahl
der Körner exponentiell. Auf dem 50. Feld wäre die Anzahl der Körner so
gross, dass sie den 365 Meter hohen Berliner Fernsehturm am
Alexanderplatz bedecken könnte. Diese Geschichte verdeutlicht die enorme
Kraft des exponentiellen Wachstums und seine beeindruckenden Ergebnisse.

Das Verständnis des Konzepts exponentiellen Wachstums ist entscheidend,
um Phänomene wie Bevölkerungswachstum, technologischen Fortschritt,
Umweltveränderungen und die Ausbreitung von Krankheiten zu analysieren.
Es macht deutlich, wie selbst kleine Veränderungen oder Entwicklungen in
einem System erhebliche Auswirkungen haben können. Die Geschichte
erinnert uns daran, dass wir sorgfältig darüber nachdenken müssen, wie
wir mit solchem Wachstum umgehen und welche langfristigen Folgen es
haben kann.

\end{tcolorbox}

Die Ressourcengewinnung hat in den letzten Jahrzehnten erheblich
zugenommen -- von 22 Milliarden Tonnen im Jahr 1970 auf 70 Milliarden
Tonnen im Jahr 2010 (UNEP, 2016). Die globale Erwärmung ist ein weiteres
drängendes Problem. Der neueste Bericht des IPCC (2023) prognostiziert
einen Anstieg der globalen Durchschnittstemperatur zwischen 1,5 und 5,8
Grad Celsius, abhängig von Szenario und künftigen Emissionen. Ein
weiteres alarmierendes Phänomen ist die Entwaldung: Alle zwei Sekunden
wird eine Waldfläche in der Grösse eines Fussballfeldes abgeholzt -- das
entspricht einer Fläche in der Grösse von New York City pro Tag. Hinzu
kommt der Rückgang der Biodiversität, der sich im Aussterben von Tier-
und Pflanzenarten zeigt und die ökologische Vielfalt unseres Planeten
bedroht.

Die grosse Beschleunigung (The Great Acceleration, Abbildungen 6 und 7)
beschreibt die beobachtbaren und messbaren (negativen) Entwicklungen
zentraler sozioökonomischer und biophysischer Indikatoren seit 1950. Die
Ursachen dieser Fehlentwicklungen können weit entfernt von dem Ort
liegen, an dem die Probleme entstehen oder ihre gravierendsten Folgen
sichtbar werden. So kann der Verlust von Arbeitsplätzen in einem Land
durch die Entscheidung eines multinationalen Unternehmens verursacht
sein, Teile seiner Tätigkeit in ein anderes Land zu verlagern, um
Profite zu maximieren. Wir können solche Probleme nicht beheben, wenn
wir nicht verstehen, wie das System -- in diesem Fall die globalisierte
Wirtschaft -- funktioniert. Ähnlich benötigen wir ein Verständnis der
globalen Klimasysteme der Erde, um uns vorzustellen, welche Folgen die
globale Erwärmung in einem bestimmten Gebiet oder einer bestimmten
Region haben könnte.

Aus diesem Grund erfordern viele Methoden und Konzepte der
Nachhaltigkeitswissenschaft systemisches Denken. Ein Ansatz systemischen
Denkens beginnt häufig mit einem Brainstorming, um ein „reichhaltiges
Bild`` aller Faktoren zu erstellen, die berücksichtigt werden müssen, um
das aktuelle Verhalten des betreffenden Systems zu verstehen. Im Fall
der Schliessung eines bestimmten Geschäftsbereichs durch ein
multinationales Unternehmen wären dies die Faktoren, die die
Entscheidung zur Fortführung, Erweiterung oder Schliessung bestimmter
Geschäftsbereiche beeinflussen könnten. Während eine erste Darstellung
möglicherweise überaus komplex oder unübersichtlich wirkt, kann die
Erstellung eines linearen Flussdiagramms helfen, die Ströme von Inputs
und Einflüssen zu verdeutlichen. Eine solche Abbildung kann bereits
erste Möglichkeiten aufzeigen, die relevanten Einflüsse neu zu bewerten
und das System so zu gestalten, dass unerwünschte Ergebnisse vermieden
werden. In manchen Fällen ist jedoch zusätzliche Arbeit notwendig, um
„konzeptionelle Modelle`` des Systems zu entwickeln, die bislang
unerkannte Verbesserungsmöglichkeiten aufzeigen.

Die Integration systemischen Denkens in die Methoden und Konzepte der
Nachhaltigkeitswissenschaft ermöglicht es uns, die Komplexität von
Problemen zu bewältigen und fundierte Strategien für eine nachhaltige
Entwicklung zu entwickeln. Systemisches Denken bildet somit eine
wichtige Grundlage, um die verschiedenen Verständnisse und Konzepte
nachhaltiger Entwicklung zu erfassen und zu gestalten, wie im nächsten
Kapitel näher ausgeführt wird.

\subsection{Nachhaltigkeit: Wann haben wir begonnen, darüber zu
sprechen?}\label{nachhaltigkeit-wann-haben-wir-begonnen-daruxfcber-zu-sprechen}

Im Jahr 1713 zeigte der sächsische Oberberghauptmann Hans Carl von
Carlowitz in seinem Werk Sylvicultura oeconomica, oder hauswirthliche
Nachricht und Naturmässige Anweisung zur wilden Baum‑Zucht die
Notwendigkeit und Möglichkeit einer „nachhaltigen Forstwirtschaft`` auf.
Eine weitere Pionierin der Nachhaltigkeit war Herzogin Anna Amalia, die
Mutter von Herzog Carl August. Sie leitete Ende des 18.Jahrhunderts die
weltweit erste Forstreform ein, die darauf abzielte, eine
kontinuierliche und gleichmässige Versorgung mit Holz sicherzustellen.
Zu jener Zeit hatte Europa einen unstillbaren Hunger nach materia prima
-- Rohstoffen, die zum Bau von Schiffen und Häusern sowie zum Kochen und
Heizen benötigt wurden. Diese Nachfrage drohte, die Ressourcen zu
erschöpfen und ihr langfristiges Bestehen zu gefährden.

Im 18. Jahrhundert beschäftigten sich Gelehrte in Deutschland und
anderen Teilen Europas mit der Endlichkeit natürlicher Ressourcen.
Anders als von Carlowitz sprach jedoch niemand ausdrücklich von
Nachhaltigkeit. Ein wichtiger Aspekt war die Versorgung einer wachsenden
Bevölkerung mit Nahrungsmitteln. Vor der Industrialisierung wurde das
Wirtschaftswachstum weitgehend von der Natur als Produktionsfaktor
bestimmt -- etwa durch die Anzahl der Bergwerke oder Schiffe oder durch
die verfügbare landwirtschaftliche Fläche. Der britische Ökonom Thomas
Robert Malthus warnte jedoch davor, dass die Nahrungsmittelproduktion
mit dem rasanten Bevölkerungswachstum infolge der Industriellen
Revolution nicht Schritt halten könne. Malthus ging davon aus, dass die
Bevölkerung exponentiell wachsen würde, während die
Nahrungsmittelproduktion nur linear ansteigen könne. Obwohl sich
Malthus' These nicht in der dramatischen Weise erfüllte, wie er es
prognostiziert hatte, gilt sein Werk als erste systematische Abhandlung
über die Grenzen des Wachstums in einer endlichen Welt.

Eine Zeit lang gerieten Malthus' Ideen in Vergessenheit, da durch die
Industrielle Revolution und den Aufstieg des Kapitalismus unbegrenztes
Wachstum möglich schien. Erst in der Mitte des 20. Jahrhunderts tauchte
die Debatte über natürliche Grenzen erneut auf -- diesmal in einer
„neo‑malthusianischen`` Perspektive, im Kontext des 1968 gegründeten
Club of Rome, der Diskussionen über planetare Grenzen und der
Wachstumskritik. Während der Malthusianismus traditionell auf externe,
physische Wachstumsgrenzen fokussiert, schlägt der ökologische Ökonom
Giorgos Kallis einen anderen Ansatz vor. Kallis (2019) plädiert für
persönliche Mässigung und freiwillige Zurückhaltung -- eine Lebensweise,
die tief in der Romantik und der antiken griechischen Philosophie
verwurzelt ist. Er fordert eine innere Begrenzung unserer Wünsche, um
uns von der ökonomischen Annahme der Knappheit und der damit verbundenen
Fixierung auf Wachstum zu befreien. Zwar anerkennt er die Realität
externer Grenzen, hält es jedoch für problematisch, wenn diese die
ökologische Argumentation gegen grenzenloses Wachstum dominieren. Kallis
stellt daher die Frage, ob es sinnvoll ist, den Klimawandel in erster
Linie als ein Problem externer Grenzen, von Knappheit oder einer
endlichen Atmosphäre zu rahmen, die keine weiteren Emissionen mehr
aufnehmen kann.

In der malthusianischen Logik müssten wir unsere Produktion immer weiter
steigern, um die Bedürfnisse einer Welt mit Grenzen zu befriedigen.
Solange diese Denkweise fortbesteht, liegt der Fokus darauf, wie wir
diese Grenzen überschreiten und die Aufnahmefähigkeit des Klimas für
eine wachsende Bevölkerung und ihre Ansprüche erweitern können.

\begin{tcolorbox}[enhanced jigsaw, opacitybacktitle=0.6, breakable, left=2mm, colback=white, leftrule=.75mm, colbacktitle=quarto-callout-note-color!10!white, bottomrule=.15mm, toprule=.15mm, titlerule=0mm, title=\textcolor{quarto-callout-note-color}{\faInfo}\hspace{0.5em}{Weiterführende Literatur}, colframe=quarto-callout-note-color-frame, bottomtitle=1mm, opacityback=0, arc=.35mm, toptitle=1mm, rightrule=.15mm, coltitle=black]

Grober U. 2013. Von Freiburg nach Rio - Carlowitz und die Bildung des
Begriffs ``Nachhaltigkeit.'' In: Die Erfindung der Nachhaltigkeit.
München: oekom verlag.

\end{tcolorbox}

\subsubsection{Das Aufeinanderprallen von Ökonomie und
Ökologie}\label{das-aufeinanderprallen-von-uxf6konomie-und-uxf6kologie}

Ursprünglich war Nachhaltigkeit ein Prinzip der Ressourcenökonomie, das
das wirtschaftliche Ziel einer langfristigen Maximierung der
Waldbewirtschaftung mit den ökologischen Bedingungen für das Nachwachsen
der Bäume verband. Im 19. Jahrhundert entstand das Konzept des „ewigen
Waldes`` mit dem Ziel, der Übernutzung ein Ende zu setzen und eine
langfristige Versorgung mit Holz -- einem lebenswichtigen Rohstoff --
sicherzustellen. Forstwissenschaftler wie Heinrich Cotta entwickelten
mathematische Methoden zur Berechnung von Holzvorräten und Zuwächsen.

Im Laufe der Zeit verlagerte sich der Fokus jedoch hin zur reinen
Ertragslehre, die auf kurze Umtriebszeiten, ertragreiche Monokulturen
und die Maximierung des finanziellen Nutzens setzte. Plötzlich stand
nicht mehr ein stetiger hoher Holzertrag im Vordergrund, sondern der
höchstmögliche direkte Geldgewinn. Natürliche Kreisläufe wurden durch
kapitalistische Dynamiken ersetzt; der Tauschwert verdrängte den
Gebrauchswert. Dies entwertete das Leitprinzip der Nachhaltigkeit, und
es dauerte über hundert Jahre -- bis in die 1960er und 70er Jahre --,
bis die wissenschaftlichen Disziplinen Ökologie und Nachhaltigkeit
wieder an Bedeutung gewannen.

\subsubsection{Club of Rome and The Limits to
Growth}\label{club-of-rome-and-the-limits-to-growth}

\begin{quote}
``If the present growth trends in world population, industrialization,
pollution, food production, and resource depletion continue unchanged,
the limits to growth on this planet will be reached sometime within the
next one hundred years.'' Meadows et al.~(1972), p.23
\end{quote}

Im Jahr\,1972 veröffentlichte der Club of Rome den Bericht The Limits to
Growth, der ein wesentlich breiteres Verständnis von „Nachhaltigkeit``
einführte. Wissenschaftler:innen begannen, einen globalen
Gleichgewichtszustand (homeostasis) zu fordern, der nur durch
koordinierte internationale Maßnahmen möglich wäre. Sie integrierten
wirtschaftliche, ökologische und soziale Aspekte von Nachhaltigkeit und
nutzten das Modell der Dynamik komplexer Systeme (system dynamics), um
die Welt besser zu verstehen. Dabei berücksichtigten sie die
Wechselwirkungen zwischen Bevölkerungsdichte, Nahrungsressourcen,
Energie, Materialien, Kapital, Umweltzerstörung und Landnutzung.
Computersimulationen verschiedener Szenarien zeigten ähnliche
Ergebnisse: einen katastrophalen Rückgang der Weltbevölkerung und des
Lebensstandards innerhalb von 50 bis 100\,Jahren, wenn die damaligen
Trends anhalten. Das Problem der ressourcen- und emissionsintensiven
Industriegesellschaften bestehe darin, dass Wachstum nicht linear,
sondern exponentiell verläuft. Langfristig führe eine solche Entwicklung
zum Zusammenbruch. Ein ökologischer Kollaps könne nur durch eine
Kurskorrektur verhindert werden, so das Fazit des Berichts.

Der Meadows-Bericht, wie er auch genannt wird, wurde stark für seine
Prognosen und Methoden kritisiert. Einige Kritiker:innen bemängelten,
dass das Systemdynamik‑Modell zu vereinfachend sei und wichtige Faktoren
wie Technologie und Innovation nicht ausreichend berücksichtige. Andere
warfen den Autor:innen vor, menschliche Anpassungsfähigkeit und die
Möglichkeit politischer Lösungen und Veränderungen zu wenig
einzubeziehen. Wieder andere kritisierten eine malthusianische
Sichtweise und einen zu pessimistischen Ausblick auf die Zukunft. Trotz
dieser Kritik entfachte der Bericht eine wichtige Debatte über
Nachhaltigkeit und die Notwendigkeit, die Umwelt zu schützen.

Ein Jahr später veröffentlichte E.F. Schumacher Small is Beautiful:
Economics as if People Mattered (1973). Schumacher stellte die
vorherrschende Vorstellung von grenzenlosem wirtschaftlichem Wachstum
und technologischem Fortschritt infrage. Stattdessen plädierte er für
eine nachhaltige, menschenzentrierte Wirtschaft, die lokale
Gemeinschaften und die Umwelt respektiert. Er argumentierte, dass eine
dezentralisierte Wirtschaft, die auf menschlichen Werten und nicht nur
auf Profit beruht, eine bessere Zukunft für alle schaffen könne. Zudem
betonte er die Bedeutung von Bildung und kultureller Entwicklung für den
Aufbau einer nachhaltigen Gesellschaft.

\subsection{Von der Entwicklungs- zur
Nachhaltigkeitsdebatte}\label{von-der-entwicklungs--zur-nachhaltigkeitsdebatte}

Die zunehmende Luft- und Wasserverschmutzung trug dazu bei, dass
Umweltthemen in Politik und Medien stärker in den Vordergrund rückten.
\emph{Greenpeace} wurde 1971 gegründet. In den 1960er- und 1970er-Jahren
untersuchten Paul Crutzen und sein Forschungsteam die Auswirkungen von
Stickoxiden auf die Ozonschicht und prognostizierten, dass diese durch
menschengemachte FCKWs stark abgebaut würde. In der Folge wurde die
Verwendung von FCKWs in Kühlschränken und Klimaanlagen verboten.

Als Reaktion auf die wachsende Bedeutung von Umweltthemen organisierte
die UNO 1972 die erste grosse Umweltkonferenz in Stockholm. Die
\emph{United Nations Conference on the Human Environment}, wie sie
genannt wurde, führte zur Gründung des \emph{United Nations Environment
Programme} (UNEP) sowie zur Einrichtung eigenständiger Umweltministerien
in vielen Ländern.

BILD

Das Entstehen von Umweltproblemen in den 1970er- und 1980er-Jahren fiel
mit den ersten „Entwicklungskrisen`` nach dem Zweiten Weltkrieg
zusammen. Damals sprach man noch von „unterentwickelten Ländern``. Nach
dem Erfolg des Marshallplans beim Wiederaufbau Europas nach dem Zweiten
Weltkrieg wurde angenommen, dass ähnliche Programme -- etwa ein
Marshallplan für Afrika -- zu einer schnellen sozioökonomischen
Aufholentwicklung der armen Länder des Globalen Südens führen würden.
Doch die 70er- und 80er-Jahre zeigten das Gegenteil. Selbst
beispielhafte Länder wie Mexiko und Brasilien gerieten in
Schuldenkrisen, was deutlich machte, dass die Erreichung
sozioökonomischer Entwicklung und die Überwindung von Armut weitaus
komplexere Herausforderungen darstellen.

Von da an war klar, dass Fragen der Entwicklung und der Umwelt auf
zwischenstaatlicher Ebene gemeinsam betrachtet werden mussten. 1983
richteten die Vereinten Nationen eine World Commission on Environment
and Development ein, die von der norwegischen Ministerpräsidentin Gro
Harlem Brundtland geleitet wurde.

\subsubsection{Die
Brundtland-Kommission}\label{die-brundtland-kommission}

Gro Harlem Brundtland berief 22 Kommissionsmitglieder, von denen drei
Viertel aus dem Globalen Süden stammten. In dieser Zeit begann die
Sowjetunion, die sich in wirtschaftlicher Stagnation befand, ihren
politischen Kurs zu ändern. Da das Wettrüsten mit den USA zu einem
erheblichen Haushaltsdefizit beigetragen hatte, führte Michail
Gorbatschow, der neu ernannte Generalsekretär der Kommunistischen Partei
der Sowjetunion, tiefgreifende Reformen ein. Gleichzeitig begann sich
ein neues Weltbild und ein globales Bewusstsein zu entwickeln. Vor
diesem Hintergrund erhielt die Brundtland-Kommission den Auftrag, einen
Bericht über die Perspektive einer globalen, nachhaltigen und
umweltfreundlichen Entwicklung bis und über das Jahr 2000 hinaus zu
erstellen.

Der 1987 vorgelegte Bericht trug den Titel \emph{Our Common Future},
wird aber meist als \emph{Brundtland-Bericht} bezeichnet.

Der Brundtland-Bericht stellte fest, dass die globalen Umweltprobleme
vor allem auf nicht-nachhaltige Konsum- und Produktionsmuster im Norden
sowie auf extreme Armut im Süden zurückzuführen sind. Die Übernutzung
der Natur und die Erschöpfung natürlicher Ressourcen verschärften die
Ungleichheiten (von Einkommen und Vermögen), erhöhten die absolute Armut
und stellten eine Bedrohung für Frieden und Sicherheit dar. Der Bericht
strebte eine faire und gerechte Definition von Nachhaltigkeit an:

\begin{quote}
Hauff Zitat hier
\end{quote}

Diese Definition brachte eine ethische Perspektive in die
Nachhaltigkeitsdebatte ein und stellte das Prinzip der Verantwortung
sowohl gegenüber heutigen als auch zukünftigen Generationen in den
Mittelpunkt. Durch den Fokus auf menschliche Bedürfnisse nahm die
Brundtland-Kommission eine anthropozentrische Position ein.

Die Brundtland-Kommission identifizierte drei zentrale Prinzipien zur
Analyse von Problemen und zur Orientierung des Handelns: Eine globale
Perspektive, die Verknüpfung von Umwelt und Entwicklung, und das Streben
nach Gerechtigkeit. Das Konzept der Gerechtigkeit wurde weiter in zwei
unterschiedliche Perspektiven unterteilt:\\
1. \textbf{Die intergenerationelle Perspektive:} Verantwortung gegenüber
zukünftigen Generationen.\\
2. \textbf{Die intragenerationelle Perspektive:} Verantwortung gegenüber
den heute lebenden Menschen, insbesondere in Entwicklungsländern, und
die Sicherstellung von Gerechtigkeit innerhalb von Ländern.

\begin{tcolorbox}[enhanced jigsaw, opacitybacktitle=0.6, breakable, left=2mm, colback=white, leftrule=.75mm, colbacktitle=quarto-callout-note-color!10!white, bottomrule=.15mm, toprule=.15mm, titlerule=0mm, title=\textcolor{quarto-callout-note-color}{\faInfo}\hspace{0.5em}{Nachhaltigkeit und nachhaltige Entwicklung}, colframe=quarto-callout-note-color-frame, bottomtitle=1mm, opacityback=0, arc=.35mm, toptitle=1mm, rightrule=.15mm, coltitle=black]

Die Konzepte „Nachhaltigkeit`` und „nachhaltige Entwicklung``
unterscheiden sich in ihrem Schwerpunkt. Nachhaltigkeit ist statisch;
sie betont einen beständigen Zustand. Nachhaltige Entwicklung hingegen
ist der dynamische Prozess, diesen Zustand zu erreichen -- sie
impliziert Bewegung und bezieht sich auf etwas, das im Entstehen ist.

\end{tcolorbox}

\subsection{1.5. Von Rio 1992 bis heute}\label{von-rio-1992-bis-heute}

Nach dem Aufruf des Brundtland-Berichts zu internationalem Handeln
richtete sich der Fokus darauf, Forderungen und Vorschläge in
verbindliche Verträge und Konventionen zu übersetzen. Die UNO wählte
hierfür das Instrument einer Konferenz, die genau 20 Jahre nach der
Stockholmer Konferenz von 1972, der ersten globalen Umweltkonferenz,
stattfand. Die United Nations Conference on Environment and Development
in Rio de Janeiro, Brasilien, im Juni 1992 -- auch bekannt als Earth
Summit -- war die bis dahin grösste internationale Konferenz mit
Delegierten aus über 170 Nationen.

Der Rio Earth Summit verabschiedete die folgenden fünf Dokumente:

\begin{itemize}
\tightlist
\item
  Eine Walddeklaration, die auf das ökologische Management und den
  Schutz der weltweiten Wälder abzielt.\\
\item
  Eine Klimaschutzkonvention, in der sich die Staaten verpflichten, die
  weltweiten Treibhausgasemissionen auf das Niveau von 1990 zu senken.\\
\item
  Eine Biodiversitätskonvention, um den Rückgang der biologischen
  Vielfalt durch völkerrechtlich verbindliche Massnahmen zu bekämpfen.\\
\item
  Die Rio-Deklaration über Umwelt und Entwicklung.\\
\item
  Agenda 21, das bekannteste der fünf Abkommen, demzufolge die
  nationalen Regierungen dafür verantwortlich sind, das
  Nachhaltigkeitsmodell in ihren Ländern umzusetzen.
\end{itemize}

Die Rio-Deklaration über Umwelt und Entwicklung betont, dass
langfristiger wirtschaftlicher Fortschritt nur durch einen
Ökosystemansatz möglich ist. Dies erfordert wiederum eine neue und
gerechte globale Partnerschaft zwischen Regierungen, Menschen und
wichtigen gesellschaftlichen Gruppen. Um dies zu erreichen, sind
internationale Vereinbarungen notwendig, die den Schutz der Umwelt und
des Entwicklungssystems sicherstellen. Die Rio-Deklaration etablierte
zentrale Prinzipien nachhaltiger Entwicklung, darunter das
Vorsorgeprinzip und das Verursacherprinzip. So besagt beispielsweise
Prinzip 15:

\begin{quote}
„In order to protect the environment, the precautionary approach shall
be widely applied by States according to their capabilities. Where there
are threats of serious or irreversible damage, lack of full scientific
certainty shall not be used as a reason for postponing cost-effective
measures to prevent environmental degradation`` (UN 1992).
\end{quote}

Im Dezember 1992 wurde die \emph{Commission on Sustainable Development
(CSD)} eingerichtet, um die wirksame Nachverfolgung der Ergebnisse des
Gipfels sicherzustellen.

\subsubsection{Agenda 21: Global denken -- lokal
handeln}\label{agenda-21-global-denken-lokal-handeln}

Agenda 21, die 40 Kapitel umfasst, behandelt alle zentralen
Politikbereiche und Handlungsfelder für nachhaltige Entwicklung. Im
Vorwort heisst es:

\begin{quote}
\emph{„\ldots integration of environment and development concerns and
greater attention to them will lead to the fulfilment of basic needs,
improved living standards for all, better protected and managed
ecosystems and a safer, more prosperous future. No nation can achieve
this on its own; but together we can -- in a global partnership for
sustainable development.``} (UN CSD, 1992)
\end{quote}

\subsubsection{Die Millenniums-Entwicklungsziele
(MDGs)}\label{die-millenniums-entwicklungsziele-mdgs}

Die Kernprinzipien der Agenda 21 -- die Stärkung von Frauen, der Schutz
der Umwelt und die Förderung einer gerechten und inklusiven Gesellschaft
-- legten den Grundstein für die Millennium Development Goals (MDGs).
Diese wurden im Jahr 2000 von der UNO verabschiedet und umfassten acht
konkrete Ziele, die bis 2015 erreicht werden sollten.

Obwohl nicht alle dieser Ziele vollständig umgesetzt wurden, schufen die
MDGs ein erhebliches Bewusstsein für die Notwendigkeit nachhaltiger
Entwicklung. Auf die MDGs folgten die Sustainable Development Goals
(SDGs), die 2015 von der UNO verabschiedet wurden.

\subsubsection{Rio+20}\label{rio20}

Die UN Conference on Sustainable Development, auch bekannt als Rio+20,
fand 2012 in Rio de Janeiro statt -- 20 Jahre nach dem historischen
Earth Summit von 1992, auf dem die Agenda 21 verabschiedet wurde.

Das Hauptziel von Rio+20 war es, das globale politische Bekenntnis zu
nachhaltiger Entwicklung zu erneuern. Die Teilnehmenden -- darunter
Regierungschefs sowie Vertreter:innen von NGOs und dem Privatsektor --
diskutierten eine Vielzahl von Themen, darunter Armutsbekämpfung,
Umweltschutz, nachhaltige Energie, Ernährungssicherheit und
Ressourcenmanagement.

Das wichtigste Ergebnis des Gipfels war die Verabschiedung einer
Erklärung mit dem Titel ``The Future We Want''. Diese Erklärung
erneuerte das Bekenntnis der internationalen Gemeinschaft zu
nachhaltiger Entwicklung, zu den Rio-Prinzipien und zu bisherigen
Aktionsplänen und skizzierte weitere Massnahmen zur Erreichung dieser
Ziele.

Darüber hinaus bereitete Rio+20 den Weg für einen universellen
Entwicklungsrahmen, der globale Nachhaltigkeitsziele und -prioritäten
über das Jahr 2015 hinaus definieren sollte.

\subsubsection{Die Agenda 2030 und Nachhaltigkeitsziele
(SDGs)}\label{die-agenda-2030-und-nachhaltigkeitsziele-sdgs}

\begin{quote}
``We are the first generation that can put an end to poverty and we are
the last generation that can put an end to climate change, so we
{[}must{]} address climate change.'' Ban-Ki Moon, UN Secretary General
2007--2016, on the 2030 Agenda
\end{quote}

Rio+20 legte den Grundstein für die Agenda 2030, die 2015 von der UNO
verabschiedet wurde. Die Agenda 2030 baut auf den Ergebnissen von Rio+20
auf und legt einen umfassenden Rahmen für nachhaltige Entwicklung fest.

Die Agenda 2030 basiert auf fünf zentralen Prinzipien bzw. Säulen (den 5
Ps) und führt die 17 Sustainable Development Goals (SDGs) ein, die
darauf abzielen, bis 2030 eine nachhaltige und gerechte Welt zu
schaffen.

\begin{tcolorbox}[enhanced jigsaw, opacitybacktitle=0.6, breakable, left=2mm, colback=white, leftrule=.75mm, colbacktitle=quarto-callout-note-color!10!white, bottomrule=.15mm, toprule=.15mm, titlerule=0mm, title=\textcolor{quarto-callout-note-color}{\faInfo}\hspace{0.5em}{Die 5 Ps der SDGs}, colframe=quarto-callout-note-color-frame, bottomtitle=1mm, opacityback=0, arc=.35mm, toptitle=1mm, rightrule=.15mm, coltitle=black]

Die 5 Ps stehen für fünf zentrale Prinzipien der Nachhaltigkeit:
\textbf{People}, \textbf{Planet}, \textbf{Prosperity}, \textbf{Peace}
und \textbf{Partnership}.

\begin{itemize}
\item
  \textbf{People:} Dieses P steht für das Ziel, soziale Gerechtigkeit,
  Chancengleichheit, gute Gesundheit und Bildung für alle Menschen zu
  fördern. Es zielt auch darauf ab, Armut zu beenden,
  Geschlechtergerechtigkeit zu fördern und das Wohlergehen der Menschen
  zu verbessern.
\item
  \textbf{Planet:} Ziel dieses P ist es, natürliche Ressourcen zu
  schützen und nachhaltig zu nutzen, die Biodiversität zu erhalten, das
  Klima zu schützen und die Umweltverschmutzung zu reduzieren. Insgesamt
  geht es darum, unseren Planeten zu bewahren und nachhaltige
  Umweltpraktiken zu fördern.
\item
  \textbf{Prosperity:} Dieses P bezieht sich auf Wirtschaftswachstum,
  nachhaltige Produktions- und Konsummuster sowie die Schaffung von
  Arbeitsplätzen und wirtschaftlichem Wohlstand für alle. Es zielt
  darauf ab, eine starke und nachhaltige Wirtschaft zu fördern, die
  allen Menschen zugutekommt.
\item
  \textbf{Peace:} Hier geht es darum, Frieden, Gerechtigkeit, gute
  Regierungsführung und starke Institutionen zu fördern. Es beinhaltet
  auch die Verhinderung von Konflikten, die Verringerung von Gewalt und
  den Aufbau inklusiver Gesellschaften.
\item
  \textbf{Partnership:} Dieses P betont die Bedeutung globaler
  Zusammenarbeit, Partnerschaften und Solidarität zwischen allen
  Akteur:innen. Es geht darum, gemeinsam an der Umsetzung der Agenda
  2030 zu arbeiten und Ressourcen, Erfahrungen und Wissen zu teilen.\\
\end{itemize}

\end{tcolorbox}

\subsubsection{Pariser Klimabakommen
2015}\label{pariser-klimabakommen-2015}

Das Pariser Klimaabkommen von 2015 ist ein wegweisendes internationales
Abkommen, das auf der 21. KOnferenz (COP 21) zur \emph{UN Framework
Convention on Climate Change (UNFCCC)} in Paris verabschiedet wurde. Das
Abkommen zielt darauf ab, den globalen Temperaturanstieg im Vergleich
zum vorindustriellen Niveau deutlich unter 2 Grad Celsius zu begrenzen
und Anstrengungen zu unternehmen, ihn auf 1,5 Grad Celsius zu begrenzen.

Das Pariser Abkommen basiert auf dem Prinzip der gemeinsamen, aber
unterschiedlichen Verantwortlichkeiten und der Berücksichtigung
nationaler Gegebenheiten. Es ermutigt alle Länder, Massnahmen zur
Reduzierung ihrer Treibhausgasemissionen zu ergreifen und sich
freiwillige Klimaziele zu setzen, die als Nationally Determined
Contributions (NDCs) bekannt sind.

Das Pariser Klimaabkommen von 2015 gilt für jedes Land, das es
ratifiziert hat, darunter auch die Schweiz. Als Vertragsstaat hat sich
die Schweiz verpflichtet, zur globalen Reduktion der
Treibhausgasemissionen beizutragen und die Ziele des Abkommens
umzusetzen.

\subsubsection{Fazit}\label{fazit}

Zwar wurden Ansätze der Nachhaltigkeit in Wirtschaft, Politik und
anderen Bereichen teilweise institutionalisiert und verbreitet, doch
eine grosse Zahl von Studien zeigt, dass sich die Welt insgesamt
weiterhin auf einem nicht-nachhaltigen Kurs befindet (siehe Voluntary
National Review of Switzerland 2022). Ein ähnliches Bild zeichnen
verschiedene Studien zur globalen Umweltsituation, wie der Global
Environment Outlook 5 (GEO-5) des UNEP von 2012, der Millennium
Development Goals Report (MDGs Report 2010) sowie Berichte zur Agenda
2030 und zum Klimawandel (IPCC 2022). Diese Studien verdeutlichen, dass
die internationale Gemeinschaft noch weit von einer nachhaltigen,
generationengerechten ökologischen, sozialen und wirtschaftlichen
Entwicklung entfernt ist.

So haben internationale Klimapolitiken den Anstieg der CO₂-Emissionen --
der Hauptursache des menschengemachten Klimawandels -- nicht gestoppt;
seit 1990 sind sie um rund 50\% gestiegen. Der Artenverlust beschleunigt
sich in vielen Regionen weiter. Der weltweite Kampf gegen Armut bleibt
hinter den gesteckten Zielen zurück, und die wirtschaftliche
Globalisierung der vergangenen drei Jahrzehnte hat die sozioökonomische
Ungleichheit in vielen Ländern verschärft, häufig verbunden mit
soziopolitischen und soziokulturellen Disparitäten.

\begin{tcolorbox}[enhanced jigsaw, opacitybacktitle=0.6, breakable, left=2mm, colback=white, leftrule=.75mm, colbacktitle=quarto-callout-caution-color!10!white, bottomrule=.15mm, toprule=.15mm, titlerule=0mm, title=\textcolor{quarto-callout-caution-color}{\faFire}\hspace{0.5em}{(Reflexions)übungen und Beispiele}, colframe=quarto-callout-caution-color-frame, bottomtitle=1mm, opacityback=0, arc=.35mm, toptitle=1mm, rightrule=.15mm, coltitle=black]

Im Selbstlernmodul ``FOKUS Onlinekurs \emph{Grundlagen Nachhaltiger
Entwicklung}'' finden Sie Reflexionsübungen, illustrative Beispiele
sowie Selbsttests zu den unterschiedlichen
Nachhaltigkeitsverständnissen.

\end{tcolorbox}

\subsection{Das Leitbild nachhaltiger Entwicklung als normatives
Konzept}\label{das-leitbild-nachhaltiger-entwicklung-als-normatives-konzept}

Das Konzept der Nachhaltigkeit in die Praxis umzusetzen und
Umsetzungsstrategien zu entwickeln, ist eine grosse Herausforderung für
unsere Gesellschaft. Es gibt unterschiedliche Verständnisse und Konzepte
von Nachhaltigkeit und nachhaltiger Entwicklung, ohne Einigkeit darüber,
wie Nachhaltigkeit erreicht und welche Massnahmen priorisiert werden
sollen. Einige dieser Verständnisse werden im Folgenden vorgestellt. Die
Vielfalt der Ansätze spiegelt die unterschiedlichen Perspektiven und
Interessen verschiedener Akteur:innen wider. Um diese Komplexität zu
bewältigen, müssen wirtschaftliche, soziale und ökologische Überlegungen
miteinander in Einklang gebracht werden. Dies erfordert einen
integrativen Ansatz. Die Herausforderung besteht darin, gemeinsame Ziele
zu definieren und Strategien zu entwickeln, um Nachhaltigkeit auf allen
Ebenen -- global, national, regional und lokal -- zu erreichen. Wir
brauchen eine breite gesellschaftliche Beteiligung sowie Dialog und
Zusammenarbeit zwischen Regierungen, Wirtschaft, Zivilgesellschaft und
Forschungseinrichtungen, um Lösungen zu finden und den notwendigen
Wandel voranzutreiben.

\begin{quote}
``When powerful metaphors become fashionable buzzwords, we risk that
diversity is accompanied by vagueness, i.e.~the phenomenon of a term
that has several meanings that `have so much in common that it is
difficult to separate them'\,'' (Feola 2015: 377)
\end{quote}

Das Konzept der Nachhaltigkeit hat in der Regel eine positive
Konnotation, kann jedoch aus unterschiedlichen Perspektiven betrachtet
werden. Damit es als Leitbild (im ökologischen, sozialen und
wirtschaftlichen Sinn) funktionieren kann, benötigt es klare Kriterien.
Doch auf welcher Grundlage können wir diese Kriterien definieren?

Nach Hirsch Hadorn \& Brun (2007) sollte ``Nachhaltigkeit'' „nachhaltige
Entwicklung``:

\begin{itemize}
\tightlist
\item
  Bedürfnisse erfüllen \ldots{}
\item
  \ldots{} auf gerechte Weise,
\item
  \ldots{} mit Blick auf die heute und künftig lebenden Menschen und
\item
  unter Berücksichtigung der Vielfalt von Werten sowie der Grenzen der
  Naturausbeutung.
\end{itemize}

Das Leitbild nachhaltiger Entwicklung ist daher nicht nur das Ergebnis
wissenschaftlicher Forschung, sondern in erster Linie ein normatives,
ethisch fundiertes Konzept. Es verbindet „ethische und analytische
Ideen`` und formuliert „Normen, die ausdrücken, was wünschenswert ist
und was geschehen soll`` (Renn et al.~2007: 39). Nachhaltige Entwicklung
ist somit ein sozialer Aushandlungs- und Entscheidungsprozess -- ein
Prozess des Suchens, Lernens und Erfahrens --, der von ethischen
Überlegungen geleitet wird. Entsprechend muss sich die
Nachhaltigkeitsforschung ihrer Einbettung in gesellschaftliche Prozesse
der Wahrnehmung und Bewertung stets bewusst sein.

\subsubsection{Von welchen Werten und Normen sollen wir uns leiten
lassen -- und
warum?}\label{von-welchen-werten-und-normen-sollen-wir-uns-leiten-lassen-und-warum}

Wenn wir vom Konzept nachhaltiger Entwicklung zu seiner Umsetzung
übergehen, kommt die Ethik ins Spiel. Ethik liefert Bewertungskriterien,
methodische Verfahren oder Prinzipien zur „Rechtfertigung und Kritik von
Handlungsregeln oder normativen Aussagen darüber, wie wir handeln
sollen`` (Fenner 2008, S.\,5). Sie hilft ausserdem, komplexe
Entscheidungen, wie sie in Dilemmasituationen auftreten, zu
strukturieren und zu begründen. Im Gegensatz zu Problemen mit
eindeutigen Lösungen sind Dilemmata komplexer und beinhalten Abwägungen
sowie das Austarieren unterschiedlicher Handlungsoptionen.

Ethik bietet daher Orientierung und Entscheidungsstrukturen, die einen
Rahmen schaffen, um in solchen Situationen einen geeigneten Handlungsweg
zu finden. Ihre Kernfunktion besteht folglich nicht darin, monokausale
Probleme zu lösen, sondern darin, komplexe Dilemmata zu strukturieren
und zu kategorisieren.

Vereinfacht gesagt bezeichnet „Moral`` persönliche oder
gesellschaftliche Überzeugungen darüber, was richtig und falsch ist,
während „Ethik`` die philosophische Auseinandersetzung mit diesen
Überzeugungen darstellt. Ethik untersucht die Regeln und Prinzipien, die
das menschliche Verhalten leiten „sollen``. Sie lässt sich in allgemeine
Ethik und angewandte Ethik unterteilen (Abbildung xy).

ABBILDUNG

Die allgemeine Ethik liefert uns Methoden und Konzepte, um grundlegende
moralische Probleme zu diskutieren. Sie besteht aus drei Teilbereichen:
normativer Ethik, deskriptiver Ethik und Metaethik. Die normative Ethik
konzentriert sich darauf, zu bestimmen, was moralisch richtig oder
falsch ist. Fragen danach, was ein gutes Leben ausmacht, können
beispielsweise je nach den individuellen Werten von Menschen
unterschiedliche Antworten haben.

Aus diesem Grund wird die normative Ethik weiter in teleologische und
deontologische Ansätze unterteilt. Die teleologische Ethik (abgeleitet
von telos, dem griechischen Wort für „Zweck``, „Ziel`` oder „Ende``)
bewertet Handlungen danach, welches Gute sie hervorbringen. Der
Utilitarismus ist beispielsweise eine teleologische Theorie, die
Handlungen nach dem daraus resultierenden Nutzen oder Glück bewertet.
Eine der ersten systematischen Ausarbeitungen des Utilitarismus ist
Jeremy Benthams (1748--1832) Werk An Introduction to the Principles of
Morals and Legislation (1780).

Die deontologische Ethik hingegen bewertet Handlungen anhand ihrer
Merkmale und nicht nach ihren Konsequenzen. Ihr Ursprung liegt im
griechischen deon („Pflicht``). Ein Beispiel für eine deontologische
Theorie ist die kantische Ethik, entwickelt vom deutschen Philosophen
Immanuel Kant.

Die Unterscheidung zwischen teleologischen und deontologischen Ansätzen
wird häufig am Beispiel des „Trolley-Problems`` diskutiert (siehe dazu
Sandel {[}2009; 2013{]}).

\begin{tcolorbox}[enhanced jigsaw, opacitybacktitle=0.6, breakable, left=2mm, colback=white, leftrule=.75mm, colbacktitle=quarto-callout-note-color!10!white, bottomrule=.15mm, toprule=.15mm, titlerule=0mm, title=\textcolor{quarto-callout-note-color}{\faInfo}\hspace{0.5em}{Weiterführende Literatur}, colframe=quarto-callout-note-color-frame, bottomtitle=1mm, opacityback=0, arc=.35mm, toptitle=1mm, rightrule=.15mm, coltitle=black]

Sandel M. 2009. Justice: What's the Right Thing to Do? Episode 01 ``THE
MORAL SIDE OF MURDER.'' https://www.youtube.com/watch?v=kBdfcR-8hEY.
Sandel M. 2013. Gerechtigkeit. Berlin: Ullstein.

\end{tcolorbox}

\subsubsection{Das Prinzip Verantwortung -- Hans
Jonas}\label{das-prinzip-verantwortung-hans-jonas}

Hans Jonas' Prinzip Verantwortung (1979) liefert einen wichtigen Beitrag
zur Nachhaltigkeitsdebatte, indem es unsere ethische Verantwortung
gegenüber zukünftigen Generationen und der Natur betont. Jonas' Ansatz,
der als eine Art „Ethik der Zukunft`` beschrieben wird, zielt darauf ab,
die neuen Herausforderungen einer technologischen Zivilisation zu
adressieren. Der Bezug zur Nachhaltigkeitsethik liegt in seinem
Verständnis von ethischer Verantwortung gegenüber zukünftigen
Generationen, die er als asymmetrische, nicht-reziproke Beziehung
begreift (Jonas 1987a, S.\,177). Diese Asymmetrie ergibt sich aus der
Macht eines „moralischen Subjekts`` über jemanden oder etwas, das
Fürsorge benötigt; Jonas beschreibt die Verwundbarkeit des Lebens
(ontologische Verwundbarkeit) und Verantwortung als unsere Pflicht,
andere Wesen vor Schaden zu bewahren (Jonas 1987a, S.\,391).

Jonas' Ansatz unterscheidet sich von anderen ethischen Theorien, indem
er die Existenz der Menschheit nicht als selbstverständlich voraussetzt
und stattdessen Pflichten gegenüber zukünftigen Generationen wie auch
gegenüber der Natur einschliesst. Jonas argumentiert, dass die erste
Pflicht einer Ethik der Zukunft darin bestehe, die fernen Folgen
menschlichen Handelns anzuerkennen (Jonas 1987a, S.\,64). Angesichts der
Unsicherheit über die langfristigen Folgen menschlichen Handelns schlägt
er ein Entscheidungsprinzip vor, das auf dem lateinischen Ausdruck in
dubio pro malo basiert, den er so interpretiert: „{[}\ldots{]} im
Zweifel ist der schlechteren Prognose der Vorzug vor der besseren zu
geben, da der Einsatz für das Spiel zu hoch geworden ist`` (Jonas 1987b,
S.\,67).

Jonas entwickelt damit einen neuen kategorischen Imperativ, der die
Notwendigkeit betont, sowohl die Natur als auch die Menschheit zu
bewahren. Dieser Imperativ gründet auf der Idee, dass die Natur einen
inhärenten Wert und Zweck hat:

\begin{quote}
„Handle so, dass die Wirkungen deiner Handlung verträglich sind mit der
Permanenz echten menschlichen Lebens auf Erden`` oder „Handle so, dass
die Wirkungen deiner Handlung nicht zerstörerisch sind für die künftigen
Möglichkeiten solchen Lebens`` (Jonas 1987a, S.36).
\end{quote}

Bereits in den späten 1970er-Jahren artikulierte Hans Jonas somit eine
grundlegende Unsicherheit über die Risiken und Schäden, die sich aus den
damals genutzten Technologien -- wie der Kernenergie -- für zukünftige
Generationen ergeben könnten. Entscheidungen über neue Technologien
erfordern eine kontinuierliche Bewertung von Risiken und Chancen, oft
unter Bedingungen unvollständiger Informationen und unsicherer
Zukunftsprognosen. Diese Entscheidungen betreffen nicht nur die Umwelt,
sondern werfen auch ethische Fragen über unsere Verpflichtung auf,
Nachhaltigkeit für zukünftige Generationen zu sichern. Jonas' Prinzip
Verantwortung betont daher den Zusammenhang zwischen Verantwortung und
Pflicht als zentrale Konzepte der Nachhaltigkeitsdebatte.

Nach Jonas haben gegenwärtige Generationen eine vorausschauende
Verantwortung gegenüber künftigen Generationen. Offen lässt er jedoch,
wie weitreichend diese Pflicht ist und ob zukünftigen Generationen
absolute oder vergleichende Standards zustehen.

Das Prinzip Verantwortung kann zudem im Zusammenhang mit verschiedenen
Verantwortungsbereichen betrachtet werden, die unterschiedliche ethische
Ansätze repräsentieren: Anthropozentrismus, Pathozentrismus,
Biozentrismus und Physiozentrismus.

ABBILDUNG

Anthropozentrismus betrachtet den Menschen als zentrales Wesen und
betont die Verantwortung gegenüber der menschlichen Gesellschaft und
ihrem Wohlergehen. Dies schliesst die Berücksichtigung der Bedürfnisse
und Interessen sowohl heutiger als auch zukünftiger Generationen ein.

Pathozentrismus hingegen betont die Verantwortung gegenüber
empfindungsfähigen Lebewesen und ihrem Wohlergehen. Er erkennt an, dass
nicht nur Menschen, sondern auch Tiere und andere Lebewesen ein Recht
darauf haben, frei von Leid und Schaden zu sein.

Biozentrismus wiederum erkennt den Eigenwert aller lebenden Organismen
an und konzentriert sich auf den Schutz und die Bewahrung der
biologischen Vielfalt. Er betont die Bedeutung des Respekts gegenüber
der natürlichen Umwelt und die Förderung nachhaltiger Praktiken.

Schliesslich fokussiert Physiozentrismus auf die Verantwortung gegenüber
der gesamten physischen Umwelt, einschliesslich der unbelebten Natur und
der Ökosysteme. Er erkennt den Eigenwert der Natur an und betont die
komplexen Wechselwirkungen zwischen allen Komponenten des Ökosystems.

Hans Jonas' Prinzip Verantwortung vereint all diese unterschiedlichen
Verantwortungsbereiche. Es fordert uns auf, Verantwortung sowohl für
Menschen als auch für andere Lebewesen, die Natur und die Umwelt zu
übernehmen. Dabei ist es wichtig, eine ausgewogene und umfassende
Perspektive einzunehmen, die langfristige Auswirkungen und das
Wohlergehen aller berücksichtigt.

Durch die Integration dieser verschiedenen Verantwortungsbereiche können
wir eine ganzheitliche und nachhaltige Ethik entwickeln, die den
Bedürfnissen und Interessen aller Aspekte des Lebens und der Natur in
angemessener Weise Rechnung trägt.

\begin{tcolorbox}[enhanced jigsaw, opacitybacktitle=0.6, breakable, left=2mm, colback=white, leftrule=.75mm, colbacktitle=quarto-callout-caution-color!10!white, bottomrule=.15mm, toprule=.15mm, titlerule=0mm, title=\textcolor{quarto-callout-caution-color}{\faFire}\hspace{0.5em}{(Reflexions)übungen und Beispiele}, colframe=quarto-callout-caution-color-frame, bottomtitle=1mm, opacityback=0, arc=.35mm, toptitle=1mm, rightrule=.15mm, coltitle=black]

Im Selbstlernmodul ``FOKUS Onlinekurs \emph{Grundlagen Nachhaltiger
Entwicklung}'' finden Sie Reflexionsübungen, illustrative Beispiele
sowie Selbsttests zur normativen Orientierung einer nachhaltigen
Entwicklung, Ethik und Dilemmatas.

\end{tcolorbox}

\subsection{Nachhaltige Entwicklung im Kontext von
Entwicklungsdebatten}\label{nachhaltige-entwicklung-im-kontext-von-entwicklungsdebatten}

Ist „nachhaltige Entwicklung`` eine Entwicklungstheorie? Ja und nein.
Nein, weil nachhaltige Entwicklung ein breiteres Gesellschaftsmodell
darstellt und als übergeordneter Rahmen fungiert, der verschiedene
Entwicklungstheorien umfasst. Ja, weil nachhaltige Entwicklung eine
normative Theorie ist, die in der wissenschaftlichen und
gesellschaftlichen Praxis de facto mit anderen Theorien der sozialen und
insbesondere der sozioökonomischen Entwicklung konkurriert.

Entwicklungstheorien analysieren soziale Entwicklung, entweder
retrospektiv oder prospektiv und entweder als Ganzes oder durch
spezifische Aspekte (z.B. wirtschaftliche Entwicklung). Die
retrospektive Untersuchung sozialer Entwicklung versucht, vergangene
oder aktuelle gesellschaftliche Entwicklungen zu erklären und Lehren für
zukünftige Planung und Entwicklungspolitik zu ziehen. Die prospektive
Untersuchung sozialer Entwicklung zielt darauf ab, gesellschaftliches
Entwicklungsmanagement und -politik zu informieren.

Entwicklungstheorien lassen sich in verschiedene Typen einteilen:
normative, strategische und explanative Theorien. Jeder Typ konzentriert
sich auf unterschiedliche Aspekte und Ziele der Entwicklungsforschung.

Normative Theorien befassen sich mit der Frage, wie Entwicklung aussehen
sollte und welche normativen Ziele sie erreichen soll. Sie untersuchen
Werte und Ziele, die mit Entwicklung verbunden sind, und beziehen sich
häufig auf Konzepte wie soziale Gerechtigkeit, Nachhaltigkeit oder
Armutsbekämpfung. Normative Theorien betonen die Definition von
Kriterien für gute Entwicklung und liefern einen normativen Rahmen für
politische Entscheidungen und Massnahmen. Beispiele sind die Theorie der
nachhaltigen Entwicklung, John Rawls' Theorie der sozialen Gerechtigkeit
oder Amartya Sens Capability Approach.

Strategische Theorien untersuchen, welche Strategien und Massnahmen
erforderlich sind, um gewünschte Entwicklungsergebnisse zu erzielen. Sie
konzentrieren sich auf Fragen der politischen Gestaltung,
Ressourcenallokation und Umsetzung von Massnahmen und analysieren, wie
bestimmte Ansätze -- wie die neoklassische Wachstumstheorie --
Entwicklung fördern können.

Explanative Theorien zielen darauf ab, die Ursachen und Mechanismen von
Entwicklung zu erklären. Sie analysieren Faktoren, die
Entwicklungsprozesse in Gesellschaften antreiben, und untersuchen die
Zusammenhänge zwischen verschiedenen Variablen. Dazu gehört die
Betrachtung wirtschaftlicher, sozialer, politischer oder kultureller
Faktoren, um Muster und Verbindungen zu identifizieren. Beispiele sind
die Modernisierungstheorie und die Dependenztheorie.

Hinweis: Alle Theorien basieren auf normativen Grundlagen -- also
zugrunde liegenden Werten und Überzeugungen -- und implizieren daher,
auch unbeabsichtigt, bestimmte Vorstellungen darüber, wie sich
Gesellschaft entwickeln sollte.

TABELLE

\begin{tcolorbox}[enhanced jigsaw, opacitybacktitle=0.6, breakable, left=2mm, colback=white, leftrule=.75mm, colbacktitle=quarto-callout-note-color!10!white, bottomrule=.15mm, toprule=.15mm, titlerule=0mm, title=\textcolor{quarto-callout-note-color}{\faInfo}\hspace{0.5em}{Entwicklungstheorien}, colframe=quarto-callout-note-color-frame, bottomtitle=1mm, opacityback=0, arc=.35mm, toptitle=1mm, rightrule=.15mm, coltitle=black]

\textbf{Modernisierungstheorie} (1950er--1980er-Jahre) betont den
Übergang von traditionellen zu modernen, industrialisierten
Gesellschaften als Grundlage für Entwicklung. Für Befürworter:innen der
Modernisierungstheorie entwickeln sich die Länder des Globalen Südens in
dieselbe Richtung wie die Industrieländer, jedoch in einem deutlich
langsameren Tempo. Die Modernisierungstheorie legt Wert auf
wirtschaftliches Wachstum, technologischen Fortschritt und
institutionelle Anpassungen. Sie wird oft mit Stufentheorien in
Verbindung gebracht, die in den 1950er- und 1960er-Jahren entstanden.
Stufentheorien betrachten Entwicklung als einen allmählichen Prozess,
der durch klar definierte Entwicklungsstufen oder -phasen gekennzeichnet
ist, mit einem besonderen Fokus auf Wirtschaftswachstum und
Modernisierung. Ein prominenter Vertreter ist Walt Rostow, der in seinem
Werk The Stages of Economic Growth: A Non-Communist Manifesto fünf
Entwicklungsstufen identifizierte (traditionelle Gesellschaft,
Voraussetzungen für den Take-off, Take-off, Übergang zur Reife und
Massenkonsum).

\textbf{Dependenztheorie} (1970er- und 1980er-Jahre) betont die
strukturelle Abhängigkeit und Ungleichheit zwischen entwickelten und
Entwicklungsländern. Sie argumentiert, dass Entwicklungsländer in einem
ungerechten globalen Wirtschaftssystem gefangen sind und in Abhängigkeit
von den Industrieländern verbleiben.

\textbf{Neoklassische Wachstumstheorie} (seit den 1990er-Jahren)
konzentriert sich auf den Zusammenhang zwischen Kapitalakkumulation,
technologischem Fortschritt und Wirtschaftswachstum. Sie betont freie
Märkte, Handel und Investitionen als Treiber von Entwicklung.

\textbf{Postkoloniale Theorien} (seit den 1990er-Jahren) betonen die
historische und strukturelle Ungleichheit zwischen ehemaligen
Kolonialmächten und Kolonien. Sie argumentieren, dass Entwicklung nur
durch eine kritische Auseinandersetzung mit dem Erbe des Kolonialismus
erreicht werden kann.

\textbf{Sens Capability Approach} (seit den 1990er-Jahren) hebt die
Bedeutung von Freiheit, sozialer Gerechtigkeit und menschlicher
Entwicklung hervor. Er betont die Notwendigkeit, die Möglichkeiten und
Fähigkeiten von Menschen zu verbessern, um Entwicklung zu erreichen.

\textbf{Post-Development-Theorien} (seit den 1990er-Jahren) kritisieren
lineare Entwicklungsmodelle, die westliche Fortschrittsvorstellungen
aufzwingen, und betonen stattdessen lokale Praktiken und Werte.
Bedeutende Denker:innen in diesem Feld sind Arturo Escobar, der in
seinem Werk Encountering Development: The Making and Unmaking of the
Third World kritisch über das Entwicklungsparadigma reflektiert, und
Vandana Shiva, die in Decolonizing the North eine nachhaltige und
gerechtere Entwicklungspolitik im Globalen Süden fordert, die den
Einfluss des Globalen Nordens reduziert.

\end{tcolorbox}

\subsection{Nachhaltigkeitsmodelle}\label{nachhaltigkeitsmodelle}

Es herrscht weitgehend Einigkeit darüber, dass nachhaltige Entwicklung
verschiedene Dimensionen umfasst und dass Nachhaltigkeit nur durch deren
Integration erreicht werden kann. Über die relative Bedeutung jeder
Dimension besteht jedoch weiterhin Uneinigkeit.

\subsubsection{Drei Dimensionen der
Nachhaltigkeit}\label{drei-dimensionen-der-nachhaltigkeit}

Das etablierte Drei-Säulen-Modell der Nachhaltigkeit, das im
Brundtland-Bericht von 1987 eingeführt wurde, versucht, ökologische,
ökonomische und soziale Belange miteinander in Einklang zu bringen.
Dieses Modell, das später auch als Triple-Bottom-Line-Modell bekannt
wurde, betont, dass wirtschaftliche Entwicklung profitabel sein sollte,
während ökologische und soziale Auswirkungen neutral bleiben.

Bei genauerer Betrachtung weist das Modell -- das ursprünglich als
Gebäude mit drei Säulen dargestellt wurde -- jedoch eine grosse Schwäche
auf. Während es suggerieren sollte, dass das System als Ganzes nicht
nachhaltig ist, wenn eine der Säulen schwach ist, könnte man
argumentieren, dass das Entfernen einer Säule -- oder sogar der beiden
äusseren Säulen -- das Gebäude nicht zum Einsturz bringen würde, solange
die verbleibende(n) Säule(n) stark genug sind.

Trotz seiner weiten Verbreitung und seines Erfolgs hat das Modell
Probleme mit Transparenz und Gerechtigkeit. So erfordert die Messung
wirtschaftlichen Erfolgs quantitative Indikatoren, während es schwierig
ist, vergleichbare Messgrössen für das soziale Wohlbefinden,
insbesondere emotionale Aspekte, zu finden. Einige Kritiker:innen
argumentieren, dass das Modell konventionelles ökonomisches Denken
priorisiert und wichtige soziale Aspekte vernachlässigt. Eine
Überbetonung menschlicher Wirtschaftssysteme kann uns zudem daran
hindern, unsere Abhängigkeit von nicht-menschlichen ökologischen
Systemen zu erkennen.

\subsubsection{Nachhaltigkeit als
Schnittmenge}\label{nachhaltigkeit-als-schnittmenge}

Ein Modell, das die drei Sektoren als Venn-Diagramm (überlappende Ringe
oder Kreise, auch bekannt als Schnittmengen- oder Triadenmodell)
darstellt, bietet eine Alternative zu den getrennten Säulen und betont
die Wechselwirkungen der drei Dimensionen. Dieses Modell verdeutlicht,
dass es engere Verbindungen zwischen zwei Bereichen geben kann und dass
die Grenzen fliessend sind.

\subsubsection{Das Vorrangmodell}\label{das-vorrangmodell}

Giddings et al.~(2002) stellen die Idee getrennter Sphären von
Wirtschaft, Gesellschaft und Umwelt infrage. Sie schlagen ein
\textbf{„verschachteltes`` Modell} vor, in dem die Wirtschaft als Teil
der Gesellschaft verstanden wird, die wiederum in der Umwelt eingebettet
ist. Dieses Modell soll sicherstellen, dass wirtschaftliche
Entscheidungen stets durch soziale und ökologische Faktoren begrenzt
werden.

Globale Ökosysteme bilden die Grundlage für soziale Systeme (d.h. die
Gesellschaft). Darin eingebettet ist die Wirtschaft, die den
Bedürfnissen der Gesellschaft dient.

Ein potenzieller Nachteil dieses Modells besteht jedoch darin, dass --
wenn die Wirtschaft als Ausgangspunkt aller Nachhaltigkeitsüberlegungen
gewählt wird -- das Risiko besteht, sie gegenüber sozialen und
ökologischen Belangen zu priorisieren.

\subsubsection{Nachhaltigkeit als
Vier-Dimensionen-Modell}\label{nachhaltigkeit-als-vier-dimensionen-modell}

Ein Problem des Drei-Säulen-Modells der Nachhaltigkeit ist, dass unklar
bleibt, was genau in die breite „soziale`` Sphäre fällt. Daher wurde
vorgeschlagen, eine vierte Dimension hinzuzufügen, um bestimmte Aspekte
hervorzuheben, die sonst im sozialen Bereich verborgen bleiben könnten.

So hat etwa der australische Kulturkommentator Jon Hawkes angeregt,
``kulturelle Vitalität`` als vierte Säule der Nachhaltigkeit neben
ökologischen, ökonomischen und sozialen Belangen aufzunehmen.

Kulturelle Nachhaltigkeit wird häufig als Teil der sozialen
Nachhaltigkeit betrachtet. Mitunter wird auch der Begriff
„soziokulturelle Nachhaltigkeit`` verwendet. In diesem Kontext bezieht
sich „Kultur`` auf Aspekte wie Gerechtigkeit,
Partizipationsmöglichkeiten, Bewusstsein für Nachhaltigkeit sowie
allgemeine Handlungs- und Verhaltensmuster (Murphy, 2017). Es ist
offensichtlich, dass die sozialen und kulturellen Dimensionen eng
miteinander verbunden sind. Kulturelle Prozesse beeinflussen unser
soziales Leben und unsere Sichtweise auf soziale Nachhaltigkeit, etwa
den Wert, den wir der Gleichheit als sozialem Ziel beimessen. Umgekehrt
prägen soziale Strukturen und Institutionen kulturelle Praktiken und
Bewertungen. Trotz dieser engen Verbindung können kulturelle und soziale
Nachhaltigkeit auch als eigenständige Dimensionen oder Perspektiven von
Nachhaltigkeit betrachtet werden.

Kultur besitzt kein eigenes Sustainable Development Goal (SDG). In der
Agenda 2030 wird Kultur im Zusammenhang mit „Zivilisation``,
„Diversität``, „Interkulturalität``, „kulturellem Erbe`` und
„Tourismus`` erwähnt -- und zwar im Rahmen der folgenden vier SDGs:
hochwertige Bildung (SDG 4), menschenwürdige Arbeit und
Wirtschaftswachstum (SDG 8), nachhaltige Städte und Gemeinden (SDG 11)
sowie nachhaltiger Konsum und Produktion (SDG 12) (Duxbury et al.,
2017).

Das Konzept kultureller Nachhaltigkeit kann auf verschiedene Weise
interpretiert werden: als vierte Dimension von Nachhaltigkeit, als
Kultur für Nachhaltigkeit und als Kultur als Nachhaltigkeit. Die
Interpretation von „Kultur als Nachhaltigkeit`` erfordert ein neues
Denken und ein neues Paradigma in Bezug auf Nachhaltigkeit. In dieser
Sichtweise wird Kultur selbst in Richtung Nachhaltigkeit transformiert.
Kulturelle Nachhaltigkeit bedeutet hier, darüber nachzudenken, was
bewahrt, was verändert und wie diese Veränderungen umgesetzt werden
sollen. Diese Interpretation -- Kultur als nachhaltige Entwicklung zu
begreifen -- definiert Kultur im weitesten Sinn als umfassenden
Lebensstil und als kontinuierlich sich verändernden Prozess. Unsere
gesellschaftliche Ordnung (wie Kapitalismus und Demokratie), unsere
Werte und unsere Arbeitsweisen sind allesamt kulturelle Produkte. Kultur
damit \ldots Kultur umfasst somit alle anderen Dimensionen der
Nachhaltigkeit, und ihre Transformation wird zu einem zentralen Anliegen
bzw. Paradigma für Nachhaltigkeit.

Die Definition kultureller Nachhaltigkeit enthält mehrere wichtige
Aspekte: Erstens betont sie die Bedeutung von intellektuellem Wachstum
und Verantwortung. Das bedeutet, Wissen zu erwerben, ein besseres
Verständnis für wichtige Themen zu entwickeln und an Bildungsaktivitäten
und -veranstaltungen teilzunehmen. Es bedeutet auch, nicht nur als
Individuen, sondern auch als Mitglieder verschiedener Gemeinschaften zu
handeln. Alle problematischen und nicht-nachhaltigen Strukturen und
Handlungsmodelle wurden von Menschen geschaffen und entwickelt. Wir
Menschen haben daher auch die Möglichkeit, sie zu verändern.

Kulturelle Nachhaltigkeit ist beispielsweise erforderlich für eine
„Grosse Transformation``, wie sie von Uwe Schneidewind (2018) gefordert
wird. In seinem Werk betont Schneidewind die Notwendigkeit umfassender
Veränderungen in allen Dimensionen der Gesellschaft, einschliesslich der
Kultur, um nachhaltige Entwicklung zu erreichen. Die Anerkennung
kultureller Nachhaltigkeit als zentrales Paradigma der Nachhaltigkeit
unterstützt die Vision einer umfassenden Transformation, die über
technologische und wirtschaftliche Lösungen hinausgeht und einen
grundlegenden gesellschaftlichen Kurswechsel anstrebt.

\subsection{Schwache und starke
Nachhaltigkeit}\label{schwache-und-starke-nachhaltigkeit}

Die wissenschaftliche Debatte über Nachhaltigkeit unterscheidet zwischen
\textbf{schwacher} und \textbf{starker Nachhaltigkeit} (vgl. z.\,B. Ott,
2020). Keines dieser Konzepte lässt sich leicht verwerfen, da beide auf
Annahmen beruhen, die nur schwer zu widerlegen sind. In der Praxis
gelten beide Konzepte als teilweise zutreffend, mit Belegen sowohl dafür
als auch dagegen.

Der zentrale Unterschied zwischen schwacher und starker Nachhaltigkeit
betrifft die Frage, \textbf{was wir erhalten sollen} und \textbf{in
welchem Mass verschiedene Kapitalformen gegeneinander austauschbar
sind}. Dabei geht es insbesondere um die folgenden drei Kapitalformen:

\textbf{Natürliches Kapital:} Die natürlichen Ressourcen und Ökosysteme,
die für das menschliche Wohlergehen und eine nachhaltige Entwicklung
entscheidend sind. Dazu gehören Land, Wasser, Luft, Biodiversität und
erneuerbare Energien. Natürliches Kapital ist die Grundlage vieler
wirtschaftlicher Aktivitäten und liefert essenzielle
Ökosystemdienstleistungen wie Nahrung, die Reinigung von Wasser und Luft
sowie kulturelle und ästhetische Werte.

\textbf{Produziertes Kapital:} Die materiellen Ressourcen und
Infrastrukturen, die in Produktion und wirtschaftlicher Tätigkeit
eingesetzt werden. Dazu zählen Gebäude, Maschinen, technische Geräte,
Transportmittel und physische Infrastrukturen. Produziertes Kapital ist
eng mit Wirtschaftswachstum und Produktivität verknüpft und spielt eine
wichtige Rolle in der wirtschaftlichen Entwicklung.

\textbf{Humankapital:} Das Wissen, die Fähigkeiten, der Bildungsstand,
die Gesundheit und die kreativen Fähigkeiten der Menschen in einer
Gesellschaft. Es umfasst individuelles und kollektives Wissen sowie die
Fertigkeiten, die für wirtschaftliche Produktivität und sozialen
Fortschritt notwendig sind. Humankapital ist entscheidend für
Innovation, Arbeitsproduktivität, gesellschaftliche Teilhabe und die
Fähigkeit einer Gesellschaft, Herausforderungen zu bewältigen und sich
anzupassen.

\subsubsection{Schwache Nachhaltigkeit}\label{schwache-nachhaltigkeit}

\textbf{Schwache Nachhaltigkeit} kann als die „weiche`` oder flexible
Option verstanden werden. Befürworter:innen der schwachen Nachhaltigkeit
gehen davon aus, dass eine Handlung dann nachhaltig ist, wenn sie dem
Gesamtsystem einen Vorteil bringt oder zumindest dessen Qualität nicht
mindert. Schwache Nachhaltigkeit kann daher als grundlegende Anforderung
an Nachhaltigkeit betrachtet werden, da sie darauf abzielt, ein
bestimmtes Niveau an Wohlergehen/Lebensqualität zu erhalten oder
Gleichheit sicherzustellen.

Das Konzept der schwachen Nachhaltigkeit „geht von der weitgehenden und
zumindest prinzipiell unbegrenzten {[}\ldots{]} Substituierbarkeit aller
Kapitalarten aus`` (Ott \& Döring 2004: 41) und basiert damit auf den
Prämissen der neoklassischen Ökonomie. Anhänger:innen der schwachen
Nachhaltigkeit glauben, dass Naturkapital durch andere Kapitalformen
ersetzt werden kann. Sie betonen das Konzept des „Gesamtkapitals`` und
argumentieren, dass die konkrete Zusammensetzung des von künftigen
Generationen geerbten Kapitals unerheblich sei, solange der Gesamtnutzen
und das Wohlergehen erhalten bleiben. Dies entspricht der neoklassischen
Nutzentheorie, nach der es unerheblich ist, wie Nutzen erzeugt wird.

Unterstützer:innen schwacher Nachhaltigkeit halten daher Massnahmen auch
dann für nachhaltig, wenn sie auf Kosten von Naturkapital durchgeführt
werden -- solange dieser Verlust durch eine Zunahme an Humankapital oder
produziertem Kapital ausgeglichen wird. So könnten beispielsweise die
Förderung von Rohstoffen wie Kohle oder der übermässige Abbau von Erdöl
gerechtfertigt werden. Auch Carlowitz' \emph{Sylvicultura
oeconomica}-Wald könnte im Prinzip ersetzt oder abgebaut werden, sofern
seine natürlichen und kulturellen Funktionen von künftigen Generationen
auf andere Weise erfüllt werden können. Dieses
Nachhaltigkeitsverständnis betont die Bedeutung des technologischen
Fortschritts für nachhaltige Entwicklung.

Die Bewertung von Massnahmen auf Grundlage der schwachen Nachhaltigkeit
ist jedoch anspruchsvoll, da sie mit erheblichen Komplexitäten und
zahlreichen Annahmen verbunden ist. Dies geht über grundlegende Fragen
zur Nachhaltigkeit hinaus, wie etwa „Was ist ein gutes Leben?{}`` oder
„Wie messen wir das Niveau des Wohlergehens?{}``. Eine zentrale
Herausforderung besteht darin, verschiedenen Kapitalarten --
insbesondere solchen ohne klaren Marktpreis -- einen monetären Wert
zuzuordnen. Welche Faktoren sollen dabei berücksichtigt werden? Welchen
Wert hat natürliche Schönheit? Welchen Wert hat ein seltener Vogel oder
ein Wal? Bullers \emph{The Value of a Whale} (2022) zeigt eindrücklich
die Probleme auf, die damit verbunden sind, der Natur einen Preis zu
geben.

Kritiker:innen des Konzepts der schwachen Nachhaltigkeit weisen auf
mehrere Einschränkungen hin. Erstens stellen sie die Annahme infrage,
dass natürliche Ressourcen unbegrenzt durch reproduzierbares Kapital
ersetzt werden können. Zweitens bezweifeln sie, dass eine Zunahme
materieller Güter den Verlust an Umweltqualität tatsächlich ausgleichen
kann. Schliesslich bestehen Bedenken hinsichtlich unserer Fähigkeit,
neue Ressourcen zu entwickeln, und der Gefahr, kritische Schwellenwerte
zu überschreiten. Diese Kritik ebnete den Weg für das Konzept der
\textbf{starken Nachhaltigkeit}.

\subsubsection{Starke Nachhaltigkeit}\label{starke-nachhaltigkeit}

Im Gegensatz zur schwachen Nachhaltigkeit betont die \textbf{starke
Nachhaltigkeit} den Erhalt des Naturkapitals. Befürworter:innen dieser
\textbf{ökologisch orientierten Theorie} sind weniger optimistisch, was
unsere Fähigkeit betrifft, natürliche Ressourcen durch produzierte
Ressourcen zu ersetzen. Sie gehen davon aus, dass diese Ressourcentypen
nur begrenzt austauschbar sind (Daly 1990; 1999; Ott 2001; Ott \& Döring
2004). Daher argumentieren sie, dass die einzelnen Elemente des
Naturkapitals so konstant wie möglich erhalten werden sollten, um
beispielsweise das Aussterben von Arten zu verhindern. Während
neoklassische Ökonom:innen zuversichtlich hinsichtlich der
Substituierbarkeit sind, betonen die Befürworter:innen der
\textbf{starken Nachhaltigkeit} die Bedeutung von \textbf{Prävention und
Vorsorge} anstelle von Nachsorge und Reaktion. So löst beispielsweise
die Reaktion auf das Ozonloch durch Sonnencreme, Schutzkleidung oder
medizinische Nachsorge nicht die Ursache des Problems. Ähnlich werden
technologische Ansätze wie das \textbf{Geo-Engineering} kritisiert, da
sie Probleme lediglich oberflächlich behandeln. Geo-Engineering
bezeichnet technische Eingriffe in geochemische oder biogeochemische
Kreisläufe, etwa zur Verlangsamung des Klimawandels oder der
Ozeanversauerung.

Bisher konzentrierte sich die Reduktion von Umweltverschmutzung
hauptsächlich auf \textbf{End-of-Pipe-Technologien}, wie
Rauchgasreinigungsanlagen an Fabrikschornsteinen oder Katalysatoren in
Autos. Dies führte jedoch auch dazu, dass langfristige Umweltprobleme
wie Klimawandel oder Biodiversitätsverlust hinausgezögert wurden,
solange deren Auswirkungen nicht unmittelbar sichtbar waren.

Dieses Phänomen steht im Zusammenhang mit dem Konzept der
\textbf{Externalisierung}, bei dem Umwelt- und Sozialkosten ausgelagert
werden, wie es Lessenich (2018) beschreibt. Das bedeutet, dass bei der
Preisbildung nur die wirtschaftlichen Kosten berücksichtigt werden --
teilweise, weil diese leichter zu quantifizieren sind --, während die
sozialen und ökologischen Kosten der Bereitstellung von Gütern und
Dienstleistungen ausgelassen oder externalisiert werden. Dies führt zu
Verzerrungen in den Marktpreisen.

\subsection{Nachhaltigkeitsstrategien}\label{nachhaltigkeitsstrategien}

Drei zentrale Strategien zeichnen sich für die Erreichung nachhaltiger
Entwicklung ab: \textbf{Effizienz}, \textbf{Konsistenz} und
\textbf{Suffizienz}. Die Erreichung von Nachhaltigkeitszielen erfordert
eine kluge Kombination dieser drei Ansätze.

ABBILDUNG

\subsubsection{Effizienzstrategie}\label{effizienzstrategie}

Die Effizienzstrategie konzentriert sich darauf, den Ressourceneinsatz
(d.h. Rohstoffe und Energie) bei der Herstellung von Produkten oder der
Erbringung von Dienstleistungen zu minimieren. Das bedeutet, den
Materialverbrauch (Materialintensität), den Energieverbrauch
(Energieintensität) und die Emissionen schädlicher Substanzen wie CO₂ zu
reduzieren.

Auch als „Öko-Effizienz`` bekannt, ist diese Strategie sowohl für
Unternehmen als auch für die Gesellschaft attraktiv, da sie Kosten,
Ressourcenverbrauch und Umweltbelastungen senken kann. Die Umsetzung
beginnt bei der Verbesserung von Produktionsprozessen, vor allem durch
technologische Innovationen. Befürworter:innen glauben, dass es möglich
ist, den Wohlstand zu verdoppeln und gleichzeitig den Naturverbrauch zu
halbieren (Weizsäcker et al.~1995).

Kritiker:innen der Effizienzstrategie sind weniger optimistisch und
warnen davor, ihre Wirkung auf nachhaltiges Wirtschaften zu
überschätzen. Sie verweisen auf Rebound-Effekte, die die
Effizienzgewinne teilweise oder sogar vollständig zunichtemachen können
(Paech 2012; Santarius 2014).

\begin{tcolorbox}[enhanced jigsaw, opacitybacktitle=0.6, breakable, left=2mm, colback=white, leftrule=.75mm, colbacktitle=quarto-callout-tip-color!10!white, bottomrule=.15mm, toprule=.15mm, titlerule=0mm, title=\textcolor{quarto-callout-tip-color}{\faLightbulb}\hspace{0.5em}{Tip}, colframe=quarto-callout-tip-color-frame, bottomtitle=1mm, opacityback=0, arc=.35mm, toptitle=1mm, rightrule=.15mm, coltitle=black]

\textbf{Umsetzung der Effizienzstrategie: Beispiele} - Energieeffiziente
Beleuchtung: Der Austausch herkömmlicher Glühbirnen durch
energieeffiziente LED-Lampen senkt den Stromverbrauch und reduziert die
Energienachfrage.

\begin{itemize}
\item
  Kraftstoffsparende Fahrzeuge: Die Entwicklung von
  kraftstoffeffizienten Hybrid- oder Elektrofahrzeugen kann zu einem
  geringeren Kraftstoffverbrauch und reduzierten Emissionen im
  Straßenverkehr führen.
\item
  Effiziente Gebäudetechnik: Die Installation intelligenter Heizungs-,
  Lüftungs- und Klimasysteme (HVAC) in Gebäuden hilft, den
  Energieverbrauch für Heizen und Kühlen zu senken.
\item
  Effiziente Wassernutzung: Der Einsatz von wassersparenden Armaturen
  und Systemen in Haushalten und Industrieunternehmen reduziert den
  Wasserverbrauch und minimiert Verschwendung.
\end{itemize}

\end{tcolorbox}

\subsubsection{Konsistenzstrategie}\label{konsistenzstrategie}

Während sich die Effizienzstrategie auf die Quantität konzentriert
(Reduzierung des Ressourceneinsatzes bei gleichzeitig gesteigertem
Output), strebt die Konsistenzstrategie (vom Lateinischen con = zusammen
+ sistere = stehenbleiben) danach, Natur und Technik in Einklang zu
bringen, indem Ressourcen wiederholt statt nur einmal genutzt werden.

Anstelle von fossilen Ressourcen, Produkten und Technologien setzt sie
auf Materialien, die mit natürlichen Stoffkreisläufen und Prozessen
kompatibel sind. Auch als „Öko-Effektivität`` bezeichnet, folgt dieses
Konzept dem Prinzip „cradle to cradle`` statt „cradle to grave`` und
basiert auf der Idee, dass intelligente Systeme ausschliesslich
Produkte, nicht aber Abfälle erzeugen.

Dies kann auf zwei Wegen erreicht werden:

Materialien können biologisch abbaubar sein (z.B. ein Shampoo aus
natürlichen Inhaltsstoffen).

Oder sie werden so gestaltet, dass sie als „technische Nährstoffe`` im
technischen Kreislauf verbleiben. Das bedeutet, dass ein ausgedientes
Produkt nicht im Müll landet, sondern in den nächsten Nutzungskreislauf
eintritt, etwa durch Upcycling (z.B. die Wiederverwendung eines
Computergehäuses als Regal).

Wie die Effizienzstrategie beginnt auch die Konsistenzstrategie mit der
Optimierung von Produktionsprozessen. Viele sehen in ihr ein grösseres
Potenzial als in der Effizienzstrategie, sowohl hinsichtlich der
Problemlösung als auch für tiefgreifende Veränderungen.

Manche Autor:innen sehen jedoch die Anwendbarkeit der Theorie als
begrenzt. So argumentieren Gerolf Hanke und Benjamin West (2013), dass
sie nicht auf alle Güter gleichermaßen übertragbar ist und die Umsetzung
erhebliche Investitionen in Produktionsanlagen und Logistik erfordern
würde. Zudem weisen sie darauf hin, dass Recyclingprozesse selbst mit
einem erhöhten Energieverbrauch verbunden sind -- im Einklang mit dem
Entropiegesetz: \textgreater{} ``Jeder materielle Wirtschaftsprozess hat
eine Zunahme der Entropie zufolge, d.h. grob vereinfacht, dass sich die
Elemente auf der stofflichen Ebene immer gleichmässiger verteilen, was
letztlich nur bedeutet, dass die Dinge und auch die Körper sich abnutzen
und eine erneute Konzentration einen immer höherer Energieaufwand
benötigt.'' Hanke \& West 2013

Ein besonderes Merkmal des Erdsystems ist die Sonnenenergie, die von
aussen kontinuierlich zugeführt wird und dem Anstieg der Entropie auf
der Erde entgegenwirkt. Dennoch wird für die Herstellung von
Energieerzeugungssystemen -- wie Biokraftstoffanlagen, Solarzellen,
Batterien für Elektroautos und Windkraftanlagen -- weiterhin sogenannte
\textbf{graue Energie} benötigt. Konsistenzstrategien können hier
Anreize für die Industrie schaffen, den Ressourcenverbrauch und die
Emissionen bei der Herstellung und Nutzung solcher Systeme zu senken.

\begin{tcolorbox}[enhanced jigsaw, opacitybacktitle=0.6, breakable, left=2mm, colback=white, leftrule=.75mm, colbacktitle=quarto-callout-tip-color!10!white, bottomrule=.15mm, toprule=.15mm, titlerule=0mm, title=\textcolor{quarto-callout-tip-color}{\faLightbulb}\hspace{0.5em}{Umsetzung der Konsistenzstrategie: Beispiele}, colframe=quarto-callout-tip-color-frame, bottomtitle=1mm, opacityback=0, arc=.35mm, toptitle=1mm, rightrule=.15mm, coltitle=black]

\begin{itemize}
\tightlist
\item
  Erneuerbare Energiesysteme: Der Übergang von fossilen Brennstoffen zu
  erneuerbaren Energiequellen wie Sonnenenergie, Windenergie und
  Wasserkraft sorgt dafür, dass die Energieproduktion im Einklang mit
  den natürlichen Rhythmen und Kreisläufen der Umwelt erfolgt.
\item
  Kreislauffähige Elektronik: Elektronische Geräte werden so konzipiert,
  dass sie leicht repariert, aufgerüstet und recycelt werden können, um
  den Ressourcenverbrauch und die Umweltbelastung zu minimieren.
\item
  Nachhaltige Bauweise: Beim Bau von Gebäuden werden nachhaltige
  Materialien eingesetzt, die geringe Umweltauswirkungen haben und nach
  Ende der Nutzungsdauer wiederverwendet oder recycelt werden können.
\end{itemize}

\end{tcolorbox}

\subsubsection{Suffizienzstrategie}\label{suffizienzstrategie}

Das Konzept der Suffizienz, abgeleitet vom lateinischen \emph{sufficere}
(„ausreichen``), hinterfragt unsere Annahmen darüber, wie viel wir für
ein gutes Leben wirklich brauchen. Es fördert eine Reduktion des
Ressourcen- und Energieverbrauchs, indem es die Nachfrage nach
ressourcenintensiven Gütern und Dienstleistungen verringert. Suffizienz,
die auch als \textbf{Genügsamkeit} oder \textbf{Angemessenheit}
beschrieben werden kann, übt Kritik an der ständigen Erzeugung neuer
Bedürfnisse durch Technologie und Werbung angesichts begrenzter
natürlicher Ressourcen. Sie fordert uns auf, nicht jedem neu
geschaffenen Bedürfnis nachzujagen und unsere Bedürfnisse auch ohne
Konsum zu erfüllen.

Während Effizienz- und Konsistenzstrategien vor allem die Produktion in
den Blick nehmen, richtet sich die Suffizienzstrategie auf den Konsum --
wenn auch nicht ausschliesslich. Sie kann in unterschiedlichem Mass und
auf verschiedenen Ebenen praktiziert werden: von kleinen
Verhaltensänderungen (z.B. Teilen statt Kaufen) bis hin zu
tiefgreifenden Veränderungen des Lebensstils (z.B. Verzicht auf
Flugreisen). Auch wenn sie beim Individuum ansetzt, lässt sich
Suffizienz auf weiteren Ebenen umsetzen, etwa durch Unternehmen
(suffizienzorientierte Produktgestaltung) oder Regierungen
(Suffizienzpolitiken). Damit sucht die Suffizienzstrategie nach dem
richtigen Gleichgewicht: \textbf{Wie viel brauchen wir für ein gutes
Leben -- und was brauchen wir nicht?}

Kritiker:innen sehen in der Suffizienzstrategie nur begrenztes Potenzial
und halten sie für schwer breit gesellschaftlich verankerbar.
Befürworter:innen dagegen betrachten sie als entscheidend für eine
nachhaltige Politik, insbesondere dort, wo Effizienz- und
Konsistenzstrategien an ihre Grenzen stossen. Das Konzept der
Postwachstumsökonomie greift die Prinzipien der Suffizienzstrategie auf.

\begin{tcolorbox}[enhanced jigsaw, opacitybacktitle=0.6, breakable, left=2mm, colback=white, leftrule=.75mm, colbacktitle=quarto-callout-tip-color!10!white, bottomrule=.15mm, toprule=.15mm, titlerule=0mm, title=\textcolor{quarto-callout-tip-color}{\faLightbulb}\hspace{0.5em}{Umsetzung der Suffizienzstrategie: Beispiele}, colframe=quarto-callout-tip-color-frame, bottomtitle=1mm, opacityback=0, arc=.35mm, toptitle=1mm, rightrule=.15mm, coltitle=black]

\begin{itemize}
\item
  Plattformen und Initiativen zum Teilen und gemeinschaftlichen Nutzen
  von Gegenständen, Werkzeugen oder Fahrzeugen ermöglichen eine
  effizientere Nutzung von Ressourcen und verringern die Anzahl
  hergestellter Produkte.
\item
  Eine überwiegend pflanzliche Ernährung senkt den Ressourcenverbrauch,
  insbesondere von Wasser und Land, im Vergleich zur Fleischproduktion.
\item
  Die Reduktion der Arbeitszeit kann zu einem geringeren
  Ressourcenverbrauch führen, da für die Produktion von Gütern und
  Dienstleistungen weniger Energie und Materialien benötigt werden.
\item
  Lokaler Konsum, etwa durch die Unterstützung regionaler Produzenten
  und Märkte, hilft, Transportwege und Emissionen zu reduzieren.
\end{itemize}

\end{tcolorbox}

\subsubsection{Rebound-Effekte: Die Kehrseite der
Effizienz}\label{rebound-effekte-die-kehrseite-der-effizienz}

ABBILDUNG un ZITAT

Kritiker:innen des Effizienzprinzips warnen vor dem Rebound-Effekt, auch
bekannt als Jevons-Paradoxon (Jevons 1865/1965). Dieses Phänomen tritt
auf, wenn Effizienzsteigerungen zu niedrigeren Preisen für Güter und
Dienstleistungen führen, dadurch die Nachfrage steigt und die
anfänglichen Einsparungen beim Ressourcenverbrauch wieder aufgehoben
werden.

Ein einfaches Beispiel: Werden Personenkraftwagen durch
Effizienzverbesserungen günstiger, entscheiden sich Käufer:innen beim
nächsten Autokauf möglicherweise für ein grösseres Modell. Zudem ist ein
kraftstoffeffizientes Auto im Betrieb günstiger, da es weniger
Treibstoff pro Kilometer benötigt. Dies motiviert viele Menschen,
häufiger oder längere Strecken mit dem Auto zu fahren und öffentliche
Verkehrsmittel oder das Fahrrad seltener zu nutzen. Dadurch werden die
technisch möglichen Effizienzgewinne in der Praxis oft nicht realisiert,
weil das Produkt häufiger oder intensiver genutzt wird (direkter
Rebound).

Darüber hinaus kann es zu weiteren Verhaltensänderungen kommen, die die
Umwelt belasten (indirekter Rebound). Im Beispiel könnte das durch die
Effizienz eingesparte Geld etwa für Flugreisen ausgegeben werden -- was
wiederum einen Teil der Umweltvorteile zunichtemacht.

ABBILDUNG

Um den \textbf{Rebound-Effekt} zu verstehen, fokussiert sich die
Literatur auf drei zentrale Bereiche: \textbf{1) finanzielle Faktoren,
2) sozio-psychologische Einflüsse und 3) regulatorische Effekte}.

\textbf{1) Finanzielle Faktoren} sind eine Hauptursache des
Rebound-Effekts. Wenn Effizienz zu Kosteneinsparungen führt, werden
finanzielle Mittel freigesetzt. Dies kann dazu führen, dass Menschen
entweder \textbf{mehr vom gleichen Produkt konsumieren} (direkter
Rebound-Effekt) oder in \textbf{andere Produkte investieren} (indirekter
Rebound-Effekt). Der finanzielle Rebound-Effekt kann auf mehreren
Mechanismen beruhen: dem \textbf{Einkommenseffekt}, bei dem
Effizienzsteigerungen reale Einkommensgewinne bewirken, dem
\textbf{Reinvestitionseffekt}, wenn Unternehmen Kosteneinsparungen zur
Ausweitung der Produktion oder für Investitionen in andere Produkte
nutzen, und dem \textbf{Marktpreiseffekt}, bei dem sinkende Nachfrage in
einem Sektor (z.\,B. effizientere Autos führen zu weniger
Kraftstoffbedarf) zu niedrigeren Preisen und dadurch zu steigender
Nachfrage in anderen energieintensiven Sektoren führt. Das Ausmass
dieses Effekts hängt von Faktoren wie Preiselastizität und
Konsumverhalten ab und kann sich vom individuellen bis auf
makroökonomische Ebene kumulieren.

\textbf{2) Sozio-psychologische Einflüsse} umfassen das „mentale
Buchführen`` von Konsument:innen. Wer ein effizientes Produkt kauft,
„erlaubt`` sich möglicherweise, \textbf{mehr zu konsumieren}.
Effizienzsteigerungen bei Produkten, die zuvor als umweltschädlich
galten, können so paradoxerweise zu \textbf{steigendem Konsum} führen
(bekannt als \textbf{Moral-Hazard-Effekt}). Wenn dieser Mehrkonsum nicht
bewusst-rational erfolgt, spricht man vom \textbf{Moral-Leaking-Effekt}
-- etwa wenn man nach dem Kauf eines sparsamen Autos mehr fährt, weil
der Kauf allein schon als umweltfreundlich empfunden wird. Der
\textbf{indirekte Rebound-Effekt} lässt sich ebenfalls durch solche
Mechanismen erklären: Der \textbf{Moral-Licensing-Effekt} beschreibt,
dass der Kauf eines ressourceneffizienten Produkts als Rechtfertigung
für einen \textbf{höheren Konsum in anderen Bereichen} dient, z.\,B.
häufigere Flugreisen nach dem Kauf eines sparsamen Autos.

\textbf{3) Regulatorische Effekte} entstehen, wenn staatliche
Vorschriften oder Anreize für effiziente Technologien
\textbf{unerwünschte Nebenwirkungen} haben. Wenn Menschen durch
Förderungen dazu ermutigt werden, grosse Elektrogeräte zu kaufen, kann
dies dazu führen, dass häufiger neue Geräte angeschafft werden -- selbst
wenn diese effizienter sind, hebt dies einen Teil der Energieeinsparung
auf. Zudem führen neue Effizienztechnologien oft zu \textbf{zusätzlichen
Kapazitäten und Infrastrukturen}, die neue Märkte schaffen (z.\,B.
Windkraftanlagen mit neuen Infrastrukturen und Arbeitsplätzen, aber auch
erheblichem Ressourcenverbrauch). Auch kommt es vor, dass
Konsument:innen \textbf{alte Produkte nach dem Kauf effizienter Geräte
weiter nutzen}, was den Gesamtverbrauch erhöht (Santarius 2012).

\begin{tcolorbox}[enhanced jigsaw, opacitybacktitle=0.6, breakable, left=2mm, colback=white, leftrule=.75mm, colbacktitle=quarto-callout-note-color!10!white, bottomrule=.15mm, toprule=.15mm, titlerule=0mm, title=\textcolor{quarto-callout-note-color}{\faInfo}\hspace{0.5em}{Rebound-Effekte}, colframe=quarto-callout-note-color-frame, bottomtitle=1mm, opacityback=0, arc=.35mm, toptitle=1mm, rightrule=.15mm, coltitle=black]

Der \textbf{direkte Rebound-Effekt} tritt auf, wenn verbesserte
Energieeffizienz den Konsum einer Ressource günstiger macht und die
Menschen daher mehr davon nutzen. Ein Beispiel: Ein
kraftstoffeffizienteres Auto ist im Betrieb günstiger, was dazu führen
kann, dass mehr gefahren wird und so die durch die Effizienzsteigerung
erzielten Energieeinsparungen teilweise oder vollständig aufgehoben
werden.

Der \textbf{indirekte Rebound-Effekt} entsteht, wenn Energieeinsparungen
oder Effizienzgewinne in einem Bereich dazu führen, dass die frei
gewordenen Ressourcen in einem anderen Bereich genutzt werden. Dadurch
verringern sich die insgesamt erzielten Einsparungen. Ein Beispiel: Das
durch energieeffizienteres Wohnen eingesparte Geld wird stattdessen für
häufigere Flugreisen ausgegeben.

\end{tcolorbox}

\textbf{Wie gross ist der Rebound-Effekt?}

Es ist schwierig, die genaue Grösse des Rebound-Effekts zu bestimmen:
Empirische Schätzungen variieren stark, abhängig von den verwendeten
Methoden und den berücksichtigten Effekten. Besonders herausfordernd ist
es, Rebound-Effekte klar von Wachstums- oder Strukturveränderungen
abzugrenzen. Auch die Art der angebotenen Produkte und Dienstleistungen
beeinflusst das Ausmass des Rebound-Effekts.

Ein Beispiel: Wird für den täglichen Arbeitsweg ein sparsameres
Transportmittel genutzt, etwa ein effizienteres Auto, sinken zwar die
Kosten, doch bleibt die verfügbare Zeit begrenzt. Man kann also nicht
beliebig häufiger oder länger pendeln, da ein Tag nur 24 Stunden hat.
Der Rebound-Effekt ist in diesem Fall relativ gering. Anders sieht es
bei Freizeitaktivitäten wie Flugreisen aus: Werden die Kosten durch
effizientere Flugzeuge oder günstigere Preise gesenkt, steigt der
Anreiz, häufiger zu fliegen -- beispielsweise ein zusätzliches
Wochenendtrip statt nur einer Reise pro Jahr. Da die Zeit hier weniger
einschränkend wirkt, kann der zusätzliche Konsum deutlich steigen -- und
damit auch der Rebound-Effekt.

Eine weitere wichtige Komponente ist der \textbf{Sättigungsgrad mit
Gütern und Dienstleistungen}. Beobachtungen zeigen, dass Rebound-Effekte
in \textbf{Hochlohnländern} tendenziell geringer ausfallen als in
Entwicklungsländern, wo noch ein erheblicher Bedarf an zusätzlichem
Konsum besteht.

So können direkte Rebound-Effekte bei Effizienzverbesserungen im Bereich
\textbf{Raumheizung} dazu führen, dass 10--30\% weniger Energie
eingespart werden als technisch möglich wäre (z.B. Hediger et al.,
2018). Die Rebound-Effekte im \textbf{Transportsektor} variieren noch
stärker: Laut Anderson et al.~(2019) könnten sie die potenziellen
Einsparungen bei privater Mobilität in Dänemark um 7,5\%, in Schweden um
30\% und in Deutschland um 60\% reduzieren. Studien zur
\textbf{Beleuchtung in Privathaushalten} zeigen dagegen sehr niedrige
Rebound-Effekte von unter 10\% (Sorrell, 2007). Im Durchschnitt können
die tatsächlichen Energieeinsparungen für diese Dienstleistungen also
bis zu 25\% unter dem technisch möglichen und prognostizierten Niveau
liegen. Das genaue Ausmass des Rebound-Effekts hängt jedoch stark vom
jeweiligen Kontext ab und kann durch geeignete Massnahmen reduziert
werden.

In seltenen Fällen können die Einsparungen sogar
\textbf{überkompensiert} werden -- dieser Effekt wird als
\textbf{„Backfire``} bezeichnet. Dabei handelt es sich jedoch um eine
Ausnahme, und der Backfire-Effekt ist kein reiner Rebound-Effekt, da er
mit Wachstums- und Strukturveränderungen verbunden ist. Ein gutes
Beispiel hierfür ist die \textbf{Digitalisierung}, die kurz- und
mittelfristig zu Backfire-Effekten führen kann (Peng et al., 2023).

\begin{tcolorbox}[enhanced jigsaw, opacitybacktitle=0.6, breakable, left=2mm, colback=white, leftrule=.75mm, colbacktitle=quarto-callout-caution-color!10!white, bottomrule=.15mm, toprule=.15mm, titlerule=0mm, title=\textcolor{quarto-callout-caution-color}{\faFire}\hspace{0.5em}{(Reflexions)übungen und Beispiele}, colframe=quarto-callout-caution-color-frame, bottomtitle=1mm, opacityback=0, arc=.35mm, toptitle=1mm, rightrule=.15mm, coltitle=black]

Im Selbstlernmodul ``FOKUS Onlinekurs \emph{Grundlagen Nachhaltiger
Entwicklung}'' finden Sie Reflexionsübungen, illustrative Beispiele
sowie Selbsttests zu den unterschiedlichen Nachhaltigkeitsmodellen und
Nachhaltigkeitsstrategien.

\end{tcolorbox}




\end{document}
